\documentclass[11pt,reqno]{amsart}

\usepackage[utf8]{inputenc}
\usepackage[T1]{fontenc}
\usepackage{amsmath,amssymb,amsthm,mathtools}
\usepackage[margin=1in]{geometry}
\usepackage{tikz-cd}
\usepackage[colorlinks=true,linkcolor=blue,citecolor=blue]{hyperref}
\usepackage{cleveref}

%---- theorem environments ---------------------------------------------
\theoremstyle{plain}
\newtheorem{theorem}{Theorem}[section]
\newtheorem{proposition}[theorem]{Proposition}
\newtheorem{lemma}[theorem]{Lemma}
\newtheorem{corollary}[theorem]{Corollary}
\newtheorem{claim}[theorem]{Claim}

\theoremstyle{definition}
\newtheorem{definition}[theorem]{Definition}
\newtheorem{remark}[theorem]{Remark}
\newtheorem{example}[theorem]{Example}
\newtheorem{setup}[theorem]{Set-up}

%---- macros -----------------------------------------------------------
\newcommand{\PP}{\mathbb{P}}
\newcommand{\PM}{\mathbb{P}^M}
\newcommand{\OO}{\mathcal{O}}
\newcommand{\sHom}{\mathcal{H}om}
\newcommand{\Hom}{\operatorname{Hom}}
\newcommand{\Ext}{\operatorname{Ext}}
\newcommand{\coker}{\operatorname{coker}}
\renewcommand{\Im}{\operatorname{Im}}
\newcommand{\id}{\mathrm{id}}
\newcommand{\Yt}{\widetilde{Y}}
\newcommand{\Ht}{\widetilde{H}}
\newcommand{\Ct}{\widetilde{C}}
\newcommand{\yt}{\tilde{y}}
\newcommand{\NY}{N_{Y}}
\newcommand{\NYt}{N_{\Yt}}
\newcommand{\Hc}{\mathcal{H}}
\newcommand{\TP}{T_{\PM}|_Y}
\newcommand{\TPt}{T_{\PM}|_{\Yt}}
\newcommand{\SC}{\mathcal{S}_C}
\newcommand{\RC}{\mathcal{R}_C}
\newcommand{\GH}{G_H}
\newcommand{\Dd}{D}
\newcommand{\cY}{\mathcal{Y}}
\newcommand{\cYt}{\widetilde{\mathcal{Y}}}
\newcommand{\cCt}{\widetilde{\mathcal{C}}}
\newcommand{\cN}{\mathcal{N}}
\newcommand{\pr}{\mathrm{pr}}

%---- title ------------------------------------------------------------
\title{Prime K3 surfaces}
\author{Jayan Mukherjee}
\address{Department of Mathematics, Oklahoma State University}
\email{jayan.mukherjee@okstate.edu}
\date{\today}

\begin{document}

\begin{abstract}
%% abstract to be written
\end{abstract}

\maketitle

\section{Introduction}

%% content to be added

\section{Preliminaries on degeneration of $K3$ surfaces into $K3$ double structures on Hirzebruch surfaces}

We recall some definitions and a theorem on some special non-reduced degenerations of $K3$ surfaces in a projective space.

\begin{definition}\label{ribbon}
Let $Y$ be a reduced connected projective variety.
 \begin{enumerate}
     \item A scheme $\widetilde{Y}$ is called a ribbon on $Y$ with conormal bundle $L$ if $\widetilde{Y}_{\operatorname{red}} = Y$, $\mathcal{I}_{Y/\widetilde{Y}}^2 = (0)$ and $\mathcal{I}_{Y/\widetilde{Y}} \cong L$ as an $\mathcal{O}_Y$ module where $\mathcal{I}_{Y/\widetilde{Y}}$ is the ideal sheaf of $Y$ inside $\widetilde{Y}$ and $L$ is a line bundle on $Y$. 

     \item If $Y$ is a smooth surface, then a ribbon $\widetilde{Y}$ on $Y$ is called a $K3$ carpet if the dualizing sheaf $K_{\widetilde{Y}} \cong \mathcal{O}_{\widetilde{Y}}$ and $H^1(\mathcal{O}_{\widetilde{Y}})= 0$.

\end{enumerate}
\end{definition}  

The following lemma characterizes the conormal bundle of a $K3$ carpet $\widetilde{Y}$ on a smooth surface $Y$.

\begin{lemma}(see \cite[\textcolor{black}{Proposition $1.5$}]{GGP08})
    A ribbon $\widetilde{Y}$ on a smooth \textcolor{black}{regular} surface $Y$ with conormal bundle $L$ is a $K3$ carpet if and only if $L \cong K_Y$.
\end{lemma}

\color{black}
We recall the definition of the Hilbert scheme of $K3$ surfaces.

\begin{definition}\label{Hilbert schemes of K3 surfaces}
    Let $\mathcal{H}_{r,g}$ denote the Hilbert scheme of $K3$ surfaces of index $r$ and genus $g$, i.e, $K3$ surfaces $X \subset \mathbb{P}^{1+r^2(g-1)} = \mathbb{P}^M$ embedded by the complete linear series of $rB$ where $B$ is a primitive very ample line bundle with $B^2 = 2g-2$. If $r = 1$, we call $\mathcal{H}_{r,g}$ the Hilbert scheme of prime $K3$ surfaces while if $r \geq 2$, $\mathcal{H}_{r,g}$ is called the Hilbert scheme of non-prime $K3$ surfaces. Recall that $\operatorname{dim}\mathcal{H}_{r,g} = h^0(N_{X/\mathbb{P}^M}) =  18+(M+1)^2$.
\end{definition}
\color{black}

Next we recall a degeneration theorem of $K3$ surfaces into $K3$ carpets on Hirzebruch surfaces.

\begin{theorem}\label{degeneration into $K3$ carpets}(\cite[ \textcolor{black}{Theorems $3.1, 4.3, 4.4, 5.1$}]{BMR21} and \cite[\textcolor{black}{Proposition $1.7$, Theorems $3.5, 4.1$}]{GP97})
Let $Y = \mathbb{F}_e$ be a Hirzebruch surface. Let $H = aC_0+bf$ be a very ample line bundle on $Y$, where $C_0$ is the class of section, satisfying $C_0^2 = -e$ and $f$ is the class of a fiber satisfying $f^2 = 0$ (this happens if and only if $b \geq ae+1$). Consider the embedding $Y \hookrightarrow \mathbb{P}^N$ induced by the complete linear series $|H|$.  
\begin{enumerate}
    \item Then there exists $K3$ carpets $\widetilde{Y}$ on $Y$ along with a very ample line bundle $\widetilde{H}$ on $\widetilde{Y}$ such that $\widetilde{H}|_Y = H$.

    \item If $\widetilde{Y} \hookrightarrow \mathbb{P}^M$ denote the embedding induced by the complete linear series $|\widetilde{H}|$, then $\widetilde{Y}$ is smoothable in $\mathbb{P}^M$ into smooth $K3$ surfaces embedded by the complete linear series of a very ample line bundle.

    \item For $0 \leq e \leq 2$, the Hilbert point of $\widetilde{Y}$ is smooth and \textcolor{black}{for} \textcolor{black}{$0 \leq e \leq 1$,} the general smooth surface of the unique Hilbert component containing $\widetilde{Y}$ is a 
    \begin{itemize}
        \item[(a)] prime $K3$ surface if $\textrm{gcd}(a,b) = 1$ 
        \item[(b)] a non-prime $K3$ surface embedded by $rB$, with $r = \operatorname{gcd}(a,b)$ and $B$ is a primitive very ample line bundle. 
    \end{itemize}
\end{enumerate}
\end{theorem}

\begin{proof} Part $(1)$, $(2)$ and the statement on smoothness of Hilbert point in part $(3)$ of the Theorem follows from \cite[Theorem $3.1 (3)$, Theorem $4.3$, Theorem $4.4$, Theorem $5.1$]{BMR21} and \cite[Proposition $1.7$, Theorem $3.5$, Theorem $4.1$]{GP97}.
    

We prove part $(3)$. Let $Y = \mathbb{F}_e$ with $0 \leq e \leq 2$. Consider the $K3$ double cover of $\pi:X \to Y$ branched along \textcolor{black}{a very general smooth} $D \in |-2K_Y|$. Let $H_0 = \pi^*(\mathcal{O}_{Y}(1)) = \pi^*(aC_0+bf)$ and consider the composition  
\[
\begin{tikzcd}
X \arrow[dr, "\varphi"] \arrow[d, "\pi"] & \\
Y \arrow[r, hook] & \mathbb{P}^M
\end{tikzcd}
\]
Recall the smoothing fibers $X_t \xhookrightarrow{\varphi_t} \mathbb{P}^M$ are obtained as the image of an embedding $\varphi_t$ which is a general deformation of $\varphi$ along a one parameter family $T$ and are hence embedded inside $\mathbb{P}^M$ by the complete linear series of a line bundle $H_t$ which is a deformation of $H_0$. Since the multiplicity of a polarization remains constant on a smooth family, we have to find the multiplicity of $H_0$ on the double cover $X$. Now since $D$ is \textcolor{black}{very general,} smooth and irreducible part $(3)$ for $Y = \mathbb{F}_e$, \textcolor{black}{$0 \leq e \leq 1$} follows from the fact that $\textrm{Pic}(X)$ is generated by $\pi^*\mathcal{O}_Y(C_0)$ and $\pi^*\mathcal{O}_Y(f)$ (see \cite[Proposition $6.3$]{AHL10} \textcolor{black}{and \cite[Proposition $6.5$ (i), (iii)]{AHL10}}).

\end{proof}

\section{Extendability of K3 carpets on rational normal scrolls}\label{sec:carpets}



\subsection{Set-up and notation}\label{subsec:setup}

Throughout this section $Y=F_e$, $e\ge 0$, is a Hirzebruch surface with ruling $p\colon Y\to\PP^1$, fibre class $f$ and negative section $C_0$ ($C_0^2=-e$), so that $K_Y=-2C_0-(e+2)f$. We fix an integer $b\ge e+3$ and put
\[
H:=C_0+bf .
\]
The complete linear series $|H|$ embeds $Y$ into $\PM$ as a rational normal scroll of type $(b-e,b)$, where
\[
M+1=h^0(\OO_Y(H))=2b-e+2 .
\]
We write $g:=M=2b-e+1$; it is the arithmetic genus of a hyperplane section of the K3 carpet below.

Let $\Yt\subset\PM$ be the K3 carpet on $Y$, i.e.\ the unique double structure on $Y$ inside $\PM$ with $I_Y/I_{\Yt}\cong K_Y$ \cite[Proposition~1.7]{GP97}. Following \cite{BE95}, $\Yt$ is determined by a surjection
\begin{equation}\label{eq:s}
s\colon \NY^{\vee}=I_Y/I_Y^2\longrightarrow K_Y ,\qquad I_{\Yt}/I_Y^2=\ker s ,
\end{equation}
and by \cite[Proposition~1.7]{GP97} such a surjection exists and is unique up to a nonzero scalar (we recover this in Lemma~\ref{lem:pushforward} below). The scheme $\Yt$ is arithmetically Cohen--Macaulay, indeed arithmetically Gorenstein \cite[Theorem~3.2]{ES18}. We put $\Ht:=\OO_{\Yt}(1)$.

\emph{Abbreviations.} We write $\NY:=N_{Y/\PM}$ and $\NYt:=N_{\Yt/\PM}$ for the normal sheaves, and
\[
\Hc:=\sHom(I_{\Yt}/I_Y^2,\OO_Y).
\]

\emph{Conventions on twists.} For a coherent sheaf $\mathcal{F}$ on $\Yt$ and a divisor class $\Dd$ on $Y$ we write $\mathcal{F}(\Dd):=\mathcal{F}\otimes_{\OO_{\Yt}}\OO_Y(\Dd)$; this is a sheaf on $Y$. Thus $\NYt(-H)=\NYt\otimes\OO_Y(-H)$ and $\NYt(-H+K_Y)=\NYt\otimes\OO_Y(-H+K_Y)$ are sheaves on $Y$, whereas $\NYt(-\Ht)=\NYt\otimes\OO_{\Yt}(-\Ht)$ is a sheaf on $\Yt$. For a sheaf on $Y$ the twist has its usual meaning. Following \cite{BM24} we set
\[
\alpha(\Yt):=h^0(\NYt(-\Ht))-(M+1),
\]
the Zak--L'vovsky invariant of $\Yt\subset\PM$.

\begin{remark}[The ribbon exact sequences]\label{rem:sequences}
For the carpet $\Yt$ we have the following exact sequences (see \cite[Lemma~4.1]{GGP08})
\begin{align}
&0\to \NYt(K_Y)\to \NYt\to \NYt\otimes\OO_Y\to 0 \label{eq:3.1}\\
&0\to \Hc\to \NYt\otimes\OO_Y\to \OO_Y(-2K_Y)\to 0 \label{eq:3.2}\\
&0\to \OO_Y(-K_Y)\to \NY\to \Hc\to 0 \label{eq:3.3}
\end{align}
In \eqref{eq:3.3} the first map is the transpose of $s$ in \eqref{eq:s} and the second is restriction of homomorphisms from $I_Y/I_Y^2$ to its subsheaf $I_{\Yt}/I_Y^2$. We shall also use the normal sequence of $Y\subset\PM$,
\begin{equation}\label{eq:3.10}
0\to T_Y\to \TP\to \NY\to 0,
\end{equation}
and the Euler exact sequence restricted to $Y$,
\begin{equation}\label{eq:3.11}
0\to \OO_Y\to \OO_Y(H)^{\oplus M+1}\to \TP\to 0 .
\end{equation}
\end{remark}

\subsection{Main theorem}

As mentioned before, to give an upper bound to $\alpha(X)$ for a $K3$ surface $X \subset \mathbb{P}^M$, we will degenerate $X \subset \mathbb{P}^M$ to a $K3$ carpet $\widetilde{Y} \subset \mathbb{P}^M$, using \Cref{degeneration into $K3$ carpets} and compute $\alpha(\widetilde{Y})$.  The main tool for all our computations that follow consist of the three exact sequences \Cref{eq:3.1}, \Cref{eq:3.2} and \Cref{eq:3.3}  that decomposes the cohomology of the twisted normal bundle $N_{\widetilde{Y}/\mathbb{P}^M}(-\widetilde{H})$ of a ribbon $\widetilde{Y} \subset \mathbb{P}^M$ in terms of cohomology of various twists of the normal bundle $N_{Y/\mathbb{P}^M}$ of the reduced part. \\

\begin{theorem}\label{thm:alpha}
Let $\Yt\subset\PM$ be the K3 carpet on the rational normal scroll $Y=F_e\subset\PM$ embedded by $|H|$, $H=C_0+bf$, and assume $b\ge e+3$. Then
\[
\alpha(\Yt)=h^0(\OO_Y(-H-2K_Y))=h^0\big(\OO_Y(3C_0+(2e+4-b)f)\big).
\]
In particular $\alpha(\Yt)=0$ if and only if $b\ge 2e+5$.
\end{theorem}

%\begin{remark}
%For $b\ge 2e+5$ this is Proposition~3.8(3) of \cite{BM25}. For $e+3\le b\le 2e+4$ (the Mukai range) it shows that the upper bound $\alpha(\Yt)\le h^0(-H-2K_Y)$ used in \cite[Theorem~4.1]{BM25} is in fact an equality.
%\end{remark}

\begin{proof}
Tensoring sequence \eqref{eq:3.1} by $-\Ht$ and taking the cohomology, we have
\begin{multline}\label{eq:3.40}
0\to H^0(\NYt(-H+K_Y))\to H^0(\NYt(-\Ht))\to H^0(\NYt(-H))\to H^1(\NYt(-H+K_Y))\\
\to H^1(\NYt(-\Ht))\to H^1(\NYt(-H))\to H^2(\NYt(-H+K_Y))\\
\to H^2(\NYt(-\Ht))\to H^2(\NYt(-H))\to 0
\end{multline}
We compute $H^i(\NYt(-H+K_Y))$ for $i=0,1,2$. Tensoring sequence \eqref{eq:3.2} by $-H+K_Y$ and taking the cohomology, we have that
\begin{multline}\label{eq:3.41}
0\to H^0(\Hc(-H+K_Y))\to H^0(\NYt(-H+K_Y))\to H^0(\OO_Y(-H-K_Y))\to H^1(\Hc(-H+K_Y))\\
\to H^1(\NYt(-H+K_Y))\to H^1(\OO_Y(-H-K_Y))\to H^2(\Hc(-H+K_Y))\to H^2(\NYt(-H+K_Y))\\
\to H^2(\OO_Y(-H-K_Y))\to 0
\end{multline}
Note that since $H=C_0+bf$, we have that $H^0(-H-K_Y)=H^0(C_0+(e+2-b)f)=0$ if $b\ge e+3$. Further, $h^2(-H-K_Y)=h^0(H+2K_Y)=h^0(-3C_0+(b-2e-4)f)=0$.

Tensoring sequence \eqref{eq:3.3} by $-H+K_Y$ and taking the cohomology, we have that
\begin{multline}\label{eq:3.42}
0\to H^0(-H)=(0)\to H^0(\NY(-H+K_Y))\to H^0(\Hc(-H+K_Y))\to H^1(-H)=(0)\\
\to H^1(\NY(-H+K_Y))\to H^1(\Hc(-H+K_Y))\to H^2(-H)=(0)\\
\to H^2(\NY(-H+K_Y))\to H^2(\Hc(-H+K_Y))\to 0
\end{multline}
Now using the Euler exact sequence \eqref{eq:3.11} tensored with $-H+K_Y$, we have that
\begin{multline}\label{eq:3.43}
0\to H^0(-H+K_Y)\to H^0(K_Y)^{\oplus M+1}=(0)\to H^0(\TP(-H+K_Y))\to H^1(-H+K_Y)\\
\to H^1(K_Y)^{\oplus M+1}\to H^1(\TP(-H+K_Y))\to H^2(-H+K_Y)\\
\to H^2(K_Y)^{\oplus M+1}\to H^2(\TP(-H+K_Y))\to 0
\end{multline}
We have $H^1(-H+K_Y)=H^1(H)^{\vee}=0$ by Kodaira vanishing, and $H^1(K_Y)=0$. Further the map $H^2(-H+K_Y)\to H^2(K_Y)^{\oplus M+1}=H^2(K_Y)\otimes H^0(H)$ is an isomorphism since it is the dual of the isomorphism $H^0(\OO_Y)\otimes H^0(H)\to H^0(H)$. Putting all this together in sequence \eqref{eq:3.43}, we have
\begin{equation}\label{eq:3.44}
H^i(\TP(-H+K_Y))=0,\qquad i=0,1,2 .
\end{equation}
On the other hand, from the sequence
\begin{equation}\label{eq:3.45}
0\to T_{Y/\PP^1}(-H+K_Y)=-C_0-(b+2)f\to T_Y(-H+K_Y)\to p^*T_{\PP^1}(-H+K_Y)=-3C_0-(b+e)f\to 0
\end{equation}
we have that, if $b\ge e+1$,
\begin{equation}\label{eq:3.46}
H^i(T_Y(-H+K_Y))=0,\ i=0,1,\qquad H^2(T_Y(-H+K_Y))=H^0(C_0+(b-2)f)^{\vee}.
\end{equation}
Hence from the sequence \eqref{eq:3.10} tensored with $-H+K_Y$ we conclude that for $b\ge e+1$,
\begin{equation}\label{eq:3.47}
\begin{gathered}
H^0(\NY(-H+K_Y))=0,\qquad H^2(\NY(-H+K_Y))=0,\\
H^1(\NY(-H+K_Y))=H^2(T_Y(-H+K_Y))=H^0(C_0+(b-2)f)^{\vee}.
\end{gathered}
\end{equation}
Plugging everything back into sequence \eqref{eq:3.42}, we have that
\begin{equation}\label{eq:3.48}
\begin{gathered}
H^0(\Hc(-H+K_Y))=0,\qquad H^2(\Hc(-H+K_Y))=0,\\
H^1(\Hc(-H+K_Y))=H^2(T_Y(-H+K_Y))=H^0(C_0+(b-2)f)^{\vee}.
\end{gathered}
\end{equation}
Going back to sequence \eqref{eq:3.41}, and noting that $H^0(-H-K_Y)=0$ for $b\ge e+3$, we have that for $b\ge e+3$,
\begin{equation}\label{eq:3.49}
H^0(\NYt(-H+K_Y))=0,\qquad H^2(\NYt(-H+K_Y))=0,
\end{equation}
and an exact sequence
\begin{equation}\label{eq:3.50}
0\to H^2(T_Y(-H+K_Y))\to H^1(\NYt(-H+K_Y))\to H^1(-H-K_Y)\to 0
\end{equation}
yielding $h^1(\NYt(-H+K_Y))=h^2(T_Y(-H+K_Y))+h^1(-H-K_Y)$.

Now we compute $H^i(\NYt(-H))$ for $i=0,1,2$. Tensoring sequence \eqref{eq:3.2} by $-H$ and taking the cohomology, we have that
\begin{multline}\label{eq:3.51}
0\to H^0(\Hc(-H))\to H^0(\NYt(-H))\to H^0(\OO_Y(-H-2K_Y))\to H^1(\Hc(-H))\to H^1(\NYt(-H))\\
\to H^1(\OO_Y(-H-2K_Y))\to H^2(\Hc(-H))\to H^2(\NYt(-H))\to H^2(\OO_Y(-H-2K_Y))\to 0
\end{multline}
Note that $h^2(-H-2K_Y)=h^0(H+3K_Y)=h^0(-5C_0+(b-3e-6)f)=0$. (In \cite{BM25} the assumption $b\ge 2e+5$, which gives $H^0(-H-2K_Y)=0$, was made at this point; we do not assume it.)

Tensoring sequence \eqref{eq:3.3} by $-H$ and taking the cohomology, we have that
\begin{multline}\label{eq:3.52}
0\to H^0(-H-K_Y)=(0)\to H^0(\NY(-H))\to H^0(\Hc(-H))\to H^1(-H-K_Y)\\
\to H^1(\NY(-H))\to H^1(\Hc(-H))\to H^2(-H-K_Y)=(0)\to H^2(\NY(-H))\to H^2(\Hc(-H))\to 0
\end{multline}
Now using the Euler exact sequence \eqref{eq:3.11} tensored with $-H$, we have that
\begin{multline}\label{eq:3.53}
0\to H^0(-H)\to H^0(\OO_Y)^{\oplus M+1}\to H^0(\TP(-H))\to H^1(-H)=(0)\to H^1(\OO_Y)^{\oplus M+1}=(0)\\
\to H^1(\TP(-H))\to H^2(-H)\to H^2(\OO_Y)^{\oplus M+1}\to H^2(\TP(-H))\to 0
\end{multline}
We have $H^0(-H)=0$, $H^1(-H)=0$ by Kodaira vanishing theorem, $H^2(-H)=0$ by pushing forward to $\PP^1$ since $a=1$. Putting all this together in sequence \eqref{eq:3.53}, we have
\begin{equation}\label{eq:3.54}
h^0(\TP(-H))=M+1,\qquad h^1(\TP(-H))=h^2(\TP(-H))=0 .
\end{equation}
On the other hand, from the sequence
\begin{equation}\label{eq:3.55}
0\to T_{Y/\PP^1}(-H)=C_0-(b-e)f\to T_Y(-H)\to p^*T_{\PP^1}(-H)=-C_0-(b-2)f\to 0
\end{equation}
we have that if $b\ge e+1$
\begin{equation}\label{eq:3.56}
H^i(T_Y(-H))=0,\ i=0,2,\qquad H^1(T_Y(-H))=H^1(C_0-(b-e)f).
\end{equation}
Hence from the sequence \eqref{eq:3.10} tensored with $-H$ we conclude that for $b\ge e+1$,
\begin{equation}\label{eq:3.57}
H^1(\NY(-H))=H^2(\NY(-H))=0
\end{equation}
and we have an exact sequence
\begin{equation}\label{eq:3.58}
0\to H^0(\TP(-H))=\mathbb{C}^{M+1}\to H^0(\NY(-H))\to H^1(T_Y(-H))\to 0 .
\end{equation}
Plugging everything back into sequence \eqref{eq:3.52}, we have that
\begin{equation}\label{eq:3.59}
H^1(\Hc(-H))=H^2(\Hc(-H))=0
\end{equation}
and an exact sequence
\begin{equation}\label{eq:3.60}
0\to H^0(\NY(-H))\to H^0(\Hc(-H))\to H^1(-H-K_Y)\to 0 .
\end{equation}
Going back to sequence \eqref{eq:3.51}, we have that for $b\ge e+3$
\begin{equation}\label{eq:3.61}
\begin{gathered}
0\to H^0(\Hc(-H))\to H^0(\NYt(-H))\to H^0(-H-2K_Y)\to 0,\\
H^1(\NYt(-H))=H^1(-H-2K_Y),\qquad H^2(\NYt(-H))=0 .
\end{gathered}
\end{equation}
Plugging \eqref{eq:3.49}, \eqref{eq:3.50}, \eqref{eq:3.61} into the exact sequence \eqref{eq:3.40}, we have
\begin{align}
h^0(\NYt(-\Ht))&=h^0(\NYt(-H))-h^1(\NYt(-H+K_Y))+h^1(\NYt(-\Ht))-h^1(\NYt(-H))\label{eq:3.62}\\
&=M+1+h^1(T_Y(-H))+h^1(-H-K_Y)+h^0(-H-2K_Y)\notag\\
&\qquad-h^2(T_Y(-H+K_Y))-h^1(-H-K_Y)+h^1(\NYt(-\Ht))-h^1(-H-2K_Y).\label{eq:3.63}
\end{align}
Now note that by \eqref{eq:3.56}, $h^1(T_Y(-H))=h^1(C_0-(b-e)f)=2b-e-2$ and $h^2(T_Y(-H+K_Y))=h^0(C_0+(b-2)f)=2b-e-2$. Hence
\begin{equation}\label{eq:3.64}
h^0(\NYt(-\Ht))=M+1+h^0(-H-2K_Y)+h^1(\NYt(-\Ht))-h^1(-H-2K_Y),
\end{equation}
that is,
\begin{equation}\label{eq:alpha-coker}
\alpha(\Yt)=h^0(-H-2K_Y)+h^1(\NYt(-\Ht))-h^1(-H-2K_Y).
\end{equation}
We end the proof by showing $h^1(\NYt(-\Ht))=h^1(-H-2K_Y)$ if $b\ge e+3$. Note that from sequence \eqref{eq:3.40} and equations \eqref{eq:3.49} and \eqref{eq:3.61}, we already have
\begin{multline}\label{eq:3.65}
0\to H^0(\NYt(-\Ht))\to H^0(\NYt(-H))\xrightarrow{\ F\ } H^1(\NYt(-H+K_Y))\\
\to H^1(\NYt(-\Ht))\to H^1(-H-2K_Y)\to 0
\end{multline}
so that $h^1(\NYt(-\Ht))=h^1(-H-2K_Y)+\dim\coker F$, and by \eqref{eq:alpha-coker}
\begin{equation}\label{eq:alpha-F}
\alpha(\Yt)=h^0(-H-2K_Y)+\dim\coker F .
\end{equation}
Hence Theorem~\ref{thm:alpha} is equivalent to the surjectivity of the connecting homomorphism $F$ in \eqref{eq:3.65}. We show that $F$ is surjective. First note that we have the sequence \eqref{eq:3.61}
\[
0\to H^0(\Hc(-H))\to H^0(\NYt(-H))\to H^0(-H-2K_Y)\to 0 .
\]
It is enough to show that the composed map
\[
H^0(\Hc(-H))\xrightarrow{\ f\ } H^1(\NYt(-H+K_Y))
\]
is surjective. Note that from \eqref{eq:3.60}, we have
\[
0\to H^0(\NY(-H))\to H^0(\Hc(-H))\xrightarrow{\ g\ } H^1(-H-K_Y)\to 0
\]
while from \eqref{eq:3.47}, \eqref{eq:3.48} and \eqref{eq:3.50}, we have the exact sequence
\[
0\to H^1(\NY(-H+K_Y))=H^2(T_Y(-H+K_Y))\to H^1(\NYt(-H+K_Y))\xrightarrow{\ h\ } H^1(-H-K_Y)\to 0 .
\]
A \v{C}ech cocycle computation (Lemma~\ref{lem:hfg} in Appendix~\ref{app:hfg}) shows that $h\circ f=2g$. Therefore we have a commutative diagram of exact sequences
\begin{equation}\label{diag:a}
\begin{tikzcd}[column sep=normal]
0\arrow[r] & H^0(\NY(-H))\arrow[r]\arrow[d,"a"'] & H^0(\Hc(-H))\arrow[r,"g"]\arrow[d,"f"] & H^1(-H-K_Y)\arrow[r]\arrow[d,"2\cdot\id"] & 0\\
0\arrow[r] & H^1(\NY(-H+K_Y))\arrow[r] & H^1(\NYt(-H+K_Y))\arrow[r,"h"] & H^1(-H-K_Y)\arrow[r] & 0
\end{tikzcd}
\end{equation}
where $a$ is the map induced by $f$ on the kernels. Since the right-hand vertical map is an isomorphism, the snake lemma gives $\coker a\cong\coker f$. Therefore, to show that $f$ is surjective, it is enough to show that $a$ is surjective. Now the proof is completed by Proposition \Cref{surjectivity of a}.
\end{proof}

\begin{proposition}\label{surjectivity of a}
The map $a: H^0(\NY(-H)) \to H^1(\NY(-H+K_Y))$ in \Cref{diag:a} is surjective.
\end{proposition}

\subsubsection*{Reformulation of the target of $a$}
We now rewrite the target $H^1(\NY(-H+K_Y))$ of $a$ by composing it with an isomorphism. The composed map will have a deformation theoretic interpretation, which will enable us to show that it is surjective. Choose a hyperplane $H\subset\PM$ (we use the same letter) meeting $Y$ transversally along a smooth curve
\[
C:=Y\cap H ,
\]
a rational normal curve of degree $M-1=g-1$ in $H\cong\PP^{g-1}$, and put $\Ct:=\Yt\cap H$, the hyperplane section of the carpet, a canonical ribbon of arithmetic genus $g$ on $C$. Since $\OO_Y(-H)=\OO_Y(-C)$, all the twists by $-H$ above are twists by $-C$. Put
\[
N_C:=N_{C/H},\qquad L_C:=K_Y|_C\cong\OO_{\PP^1}(-g-1),\qquad s_C:=s|_C\in H^0(N_C\otimes L_C),
\]
where $K_Y\cdot C=-g-1$ and $N_Y|_C\cong N_C$ by transversality.

\begin{lemma}\label{lem:pushforward}
Let $Y'\to\PP^1$ be any smooth rational normal scroll in $\PM$ of degree $M-1$ (a $\PP^1$-bundle over $\PP^1$ embedded by a line bundle of fibre degree $1$), with ruling $p$. Then
\[
p_*(N_{Y'/\PM}\otimes K_{Y'})\cong\OO_{\PP^1},\qquad R^1p_*(N_{Y'/\PM}\otimes K_{Y'})=0 .
\]
Consequently $h^0(N_{Y'/\PM}\otimes K_{Y'})=1$ and $h^1(N_{Y'/\PM}\otimes K_{Y'})=0$. Moreover, if $s'$ spans $H^0(N_{Y'/\PM}\otimes K_{Y'})$ and $\partial\colon H^0(N_{Y'/\PM}\otimes K_{Y'})\to H^1(T_{Y'}\otimes K_{Y'})$ is the connecting map of the normal sequence tensored with $K_{Y'}$, then the image of $\partial(s')$ under the Leray map
\[
H^1(T_{Y'}\otimes K_{Y'})\longrightarrow H^0\big(R^1p_*(T_{Y'}\otimes K_{Y'})\big)=H^0(\OO_{\PP^1})=k
\]
is nonzero.
\end{lemma}

\begin{proof}
Write $Y'=Y$, $H$ for the hyperplane class. From $0\to T_{Y/\PP^1}\to T_Y\to p^*T_{\PP^1}\to 0$ tensored with $K_Y$, using $T_{Y/\PP^1}\otimes K_Y=p^*K_{\PP^1}$, $p_*K_Y=0$, $R^1p_*K_Y=K_{\PP^1}$ (relative duality), we get $p_*(T_Y\otimes K_Y)=K_{\PP^1}=\OO(-2)$ and $R^1p_*(T_Y\otimes K_Y)\cong T_{\PP^1}\otimes R^1p_*K_Y=\OO_{\PP^1}$. From the Euler sequence \eqref{eq:3.11} tensored with $K_Y$, using $Rp_*(K_Y(H))=0$ (the fibre degree of $K_Y(H)$ is $-1$), we get $p_*(\TP\otimes K_Y)\cong R^1p_*K_Y=\OO(-2)$ and $R^1p_*(\TP\otimes K_Y)=0$. Pushing forward \eqref{eq:3.10} tensored with $K_Y$ gives the exact sequence
\[
0\to p_*(T_Y\otimes K_Y)\to p_*(\TP\otimes K_Y)\to p_*(\NY\otimes K_Y)\to R^1p_*(T_Y\otimes K_Y)\to R^1p_*(\TP\otimes K_Y)=0 .
\]
The first map is an injection $\OO(-2)\to\OO(-2)$, hence an isomorphism. Therefore $p_*(\NY\otimes K_Y)\xrightarrow{\ \sim\ }R^1p_*(T_Y\otimes K_Y)\cong\OO_{\PP^1}$, and $R^1p_*(\NY\otimes K_Y)=0$ because $R^1p_*(\TP\otimes K_Y)=0=R^2p_*$. Taking global sections gives $h^0=1$, $h^1=0$. Finally, the composite of $\partial$ with the Leray map is $H^0$ of the sheaf-level connecting map $p_*(\NY\otimes K_Y)\to R^1p_*(T_Y\otimes K_Y)$, which we have just shown to be an isomorphism.
\end{proof}

\begin{lemma}\label{lem:delta}
Let $\delta\colon H^0(N_C\otimes L_C)\to H^1(\NY(-H+K_Y))$ be the connecting map of the restriction sequence
\[
0\to \NY\otimes K_Y(-C)\to \NY\otimes K_Y\to (\NY\otimes K_Y)|_C=N_C\otimes L_C\to 0 .
\]
Then $\delta$ is surjective with kernel $k\cdot s_C$. Hence it induces an isomorphism
\[
\bar\delta\colon H^0(N_C\otimes L_C)/k\,s_C\ \xrightarrow{\ \sim\ }\ H^1(\NY(-H+K_Y)) .
\]
\end{lemma}

\begin{proof}
By Lemma~\ref{lem:pushforward}, $H^0(\NY\otimes K_Y)=k\,s$ and $H^1(\NY\otimes K_Y)=0$; the long exact sequence of the restriction sequence gives the claim. (As a check on dimensions: $N_C\cong\OO_{\PP^1}(g+1)^{\oplus g-2}$ for a rational normal curve of degree $g-1$, so $h^0(N_C\otimes L_C)=g-2$, while $h^1(\NY(-H+K_Y))=2b-e-2=g-3$ by \eqref{eq:3.47}.)
\end{proof}



\begin{proof}[Proof of \Cref{surjectivity of a}]
Composing $a$ with $\bar\delta^{-1}$ we conclude that it is enough to show that the map
\begin{equation}\label{eq:themap}
\bar\delta^{-1}\circ a\colon\ H^0(\NY(-H))\longrightarrow H^0(N_C\otimes K_Y|_C)\big/\,k\,s_C
\end{equation}
is surjective. Now the proof follows from \Cref{prop:a-surj}.   
\end{proof}


\subsection{Deformation theoretic interpretation of $\bar\delta^{-1}\circ a$ and proof of \Cref{prop:a-surj}}\label{subsec:geometric}

In this and the following subsections $\Dd$ denotes the scheme $\operatorname{Spec}k[\varepsilon]/(\varepsilon^2)$, and for a closed subscheme $Z\subset\PM$ we write $Z_\Dd:=Z\times\Dd$ for the constant family, a closed subscheme of $\PM_\Dd$. We use freely the identification of first-order embedded deformations of a closed subscheme $Z\subset\PM$ with sections of $N_{Z/\PM}=\sHom(I_Z,\OO_Z)$ \cite[\S3.2]{Ser06}: to a section $\eta$ corresponds the closed subscheme $Z_\varepsilon\subset\PM_\Dd$ with ideal
\[
I_{Z_\varepsilon}=\{\,f+\varepsilon g\ :\ f\in I_Z,\ g\in\OO_{\PM},\ g\equiv\eta(f)\bmod I_Z\,\}\subset\OO_{\PM}[\varepsilon]/(\varepsilon^2),
\]
the preimage of the graph of $\eta$. Locally $I_{Z_\varepsilon}$ is generated by the elements $f+\varepsilon\tilde\eta(f)$, $f\in I_Z$, where $\tilde\eta(f)$ is any local section of $\OO_{\PM}$ lifting $\eta(f)\in\OO_Z=\OO_{\PM}/I_Z$; the ideal does not depend on the choice of lifts, since $\varepsilon g=\varepsilon\,(g+\varepsilon\tilde\eta(g))$ for $g\in I_Z$.

\subsubsection*{Three parameter spaces}

%% ---- Lemma 2.6 (replaces the lemma and its proof) ----

\begin{lemma}[Scrolls through $C$]\label{lem:SC}
Let $\mathcal F$ be the flag Hilbert scheme of pairs $(C'\subset Y')$ of closed subschemes of $\PM$ with Hilbert polynomials
\[
P_{C'}(n)=(2b-e)\,n+1,\qquad P_{Y'}(n)=(2b-e)\binom{n+1}{2}+n+1,
\]
the Hilbert polynomials of $C$ and $Y$ (both of degree $2b-e=M-1$), and let $\SC$ be the fibre over the point $[C]$ of the Hilbert scheme of rational normal curves of the open subscheme of $\mathcal F$ where $Y'$ is a smooth rational normal scroll. Thus $\SC$ is the locus of smooth rational normal scrolls $Y'\subset\PM$ containing $C$; for such $Y'$, $Y'\cap H$ is an effective divisor of class $H|_{Y'}$ containing $C$ with $\deg C=H^2$, so $Y'\cap H=C$, and $H$ is the only hyperplane with this property, since $C$ spans $H$.
\begin{enumerate}
\item $T_{[Y]}\SC=H^0(\NY(-H))$. More precisely, a first-order deformation $Y_\varepsilon$ of $Y$ with class $x\in H^0(\NY)$ contains $C_\Dd$ if and only if $x\in H^0(\NY(-C))=H^0(\NY(-H))$.
\item $\SC$ is smooth at $[Y]$ (of dimension $h^0(\NY(-H))=4b-2e$), hence reduced near $[Y]$.
\item Let $\GH:=\{\gamma\in PGL_{M+1}(k):\ \gamma|_H=\id_H\}$; its Lie algebra is $H^0(T_{\PM}(-H))\cong k^{M+1}$. $\GH$ acts on $\SC$, and the differential at $1$ of the orbit map $\gamma\mapsto\gamma Y$ is the map
\[
c\colon H^0(\TP(-H))=\mathbb{C}^{M+1}\longrightarrow H^0(\NY(-H))
\]
of \eqref{eq:3.58}, which is injective. Thus $\Im c=T_{[Y]}(\GH\cdot Y)$, and the quotient $H^0(\NY(-H))/\Im c\cong H^1(T_Y(-H))$ has dimension $2b-e-2=g-3$.
\end{enumerate}
\end{lemma}

\begin{proof}
Each of the statements (1)--(3) are standard for the flag Hilbert scheme (see \cite{Kle81}, \cite[\S4.5]{Ser06} and \cite{CLM93}); For the dimension computations of the cohomology groups appearing in (1)--(3) see \eqref{eq:3.54}, \eqref{eq:3.56}, \eqref{eq:3.57} and \eqref{eq:3.58}.    
\end{proof}

%% ---- bibliography: add before \bibitem[Lvo92] ----

We recall some well-known facts about the space of canonical ribbons on a rational curve $C$ and the tangent space to that space at a given ribbon $\widetilde{C}$.

\begin{lemma}[Canonical ribbons on $C$]\label{lem:RC}
Let $\RC:=\PP\big(H^0(N_C\otimes L_C)\big)\cong\PP^{g-3}$, and recall that $s_C=s|_C\colon N_C^{\vee}\cong N_Y^{\vee}|_C\to K_Y|_C=L_C$ is the restriction to $C$ of the surjection $s$ defining the carpet, an element of $H^0(N_C\otimes L_C)$.
\begin{enumerate}
\item Every nonzero $t\in H^0(N_C\otimes L_C)=\Hom(N_C^{\vee},L_C)$ is surjective and defines a double structure $\Ct_t\subset H$ on $C$ by $I_{\Ct_t}:=\ker\big(I_C\to N_C^{\vee}\xrightarrow{t}L_C\big)$; it satisfies $I_C^2\subset I_{\Ct_t}$ and $I_C/I_{\Ct_t}\cong L_C$, and depends only on $[t]\in\RC$. Conversely every double structure on $C$ in $H$ with conormal bundle $L_C$ is of this form. 
%One has $\Ct=\Ct_{s_C}$.
\item $T_{[\Ct]}\RC=H^0(N_C\otimes L_C)/k\,s_C$. Concretely, a first-order deformation of $\Ct$ inside $H$ which is a double structure on $C_\Dd$ with invertible conormal is $\ker(s_C+\varepsilon t)$ for some $t\in H^0(N_C\otimes L_C)$, and $s_C+\varepsilon t$, $s_C+\varepsilon t'$ define the same subscheme if and only if $t-t'\in k\,s_C$.
\end{enumerate}
\end{lemma}

\begin{proof}
$(1)$ For part $(1)$ see \cite[\S1]{BE95}. \\
%(1) $N_C\cong\OO(g+1)^{\oplus g-2}$ and $L_C\cong\OO(-g-1)$, so $\Hom(N_C^{\vee},L_C)=\Hom(\OO(-g-1)^{g-2},\OO(-g-1))=k^{g-2}$ consists of constant maps, every nonzero one of which is surjective. The correspondence between double structures with conormal $L_C$ and surjections $N_C^{\vee}\to L_C$ is \cite[\S1]{BE95}. That $\Ct=\Yt\cap H$ corresponds to $s_C$: $I_{\Ct/H}=(I_{\Yt}+(h))/(h)$ and $I_{C/H}=(I_Y+(h))/(h)$, where $h$ is the equation of $H$; since $h$ is a nonzerodivisor modulo $I_Y$, $I_Y\cap(h)=hI_Y$, whence $I_{C/H}/I_{\Ct/H}=I_Y/(I_{\Yt}+hI_Y)=K_Y\otimes\OO_C=L_C$ and the surjection $N_C^{\vee}\to L_C$ is $s\otimes\OO_C=s_C$.
(2) An invertible sheaf on $C_\Dd$ restricting to $L_C$ is $L_C\otimes\OO_\Dd$ (as $H^1(\OO_C)=0$), with automorphisms $\OO_\Dd^{\times}$; so such a deformation is the kernel of a surjection $s_C+\varepsilon t\colon N_C^{\vee}\otimes\OO_\Dd\to L_C\otimes\OO_\Dd$, well defined up to multiplication by a unit $1+\varepsilon\lambda$, i.e.\ up to $t\mapsto t+\lambda s_C$.
\end{proof}

In the following lemma we prove a slight generalization of \cite{GP97}, to show that there exists a unique $K3$ carpet structure on a first order deformation $Y_{\varepsilon} \subset \mathbb{P}_D^M$ of the scroll $Y \subset \mathbb{P}^M$.

\begin{lemma}[The carpet on a deformed scroll]\label{lem:Phi}
Let $Y_\varepsilon\subset\PM_\Dd$ be a first-order deformation of $Y$. There is a unique closed subscheme $\Yt_\varepsilon\subset\PM_\Dd$, flat over $\Dd$, with $\Yt_\varepsilon\times_\Dd\{0\}=\Yt$, $Y_\varepsilon\subset\Yt_\varepsilon$, $I_{Y_\varepsilon}^2\subset I_{\Yt_\varepsilon}$ and $I_{Y_\varepsilon}/I_{\Yt_\varepsilon}$ invertible on $Y_\varepsilon$. We write $\Phi(Y_\varepsilon):=\Yt_\varepsilon$ and $d\Phi\colon H^0(\NY)\to H^0(\NYt)$ for the resulting map. Moreover:
\begin{enumerate}
\item the image of $d\Phi(x)$ under $H^0(\NYt) = \operatorname{Hom}(I_{\widetilde Y}, \mathcal O_{\widetilde Y})\to H^0(\NYt\otimes\OO_Y) = \operatorname{Hom}(I_{\widetilde Y}, \mathcal O_{Y})$ is $\iota(x):=x|_{I_{\Yt}}$, the image of $x$ under the composite $H^0(\NY) = \operatorname{Hom}(I_{Y}, \mathcal O_{Y})\to H^0(\Hc)\to H^0(\NYt\otimes\OO_Y) = \operatorname{Hom}(I_{\widetilde Y}, \mathcal O_{Y})$ of the maps in \eqref{eq:3.3} and \eqref{eq:3.2};
\item for $\gamma_\varepsilon\in PGL_{M+1}(\Dd)$ one has $\Phi(\gamma_\varepsilon Y)=\gamma_\varepsilon\Yt$.
\end{enumerate}
\end{lemma}

\begin{proof}
Such a $\Yt_\varepsilon$ is the kernel of a surjection $s_\varepsilon\colon N_{Y_\varepsilon/\PM_\Dd}^{\vee}\to L_\varepsilon$ onto an invertible sheaf with $L_\varepsilon|_Y\cong K_Y$. Since $H^1(\mathcal{O}_Y) = 0$, we have that $\operatorname{Pic}(Y_\varepsilon)\to\operatorname{Pic}(Y)$ is injective. Hence  $L_\varepsilon\cong\omega_{Y_\varepsilon/\Dd}$, and $s_\varepsilon$ is determined by $\Yt_\varepsilon$ up to $H^0(\OO_{Y_\varepsilon}^{\times})=(k[\varepsilon]/\varepsilon^2)^{\times}$. Conversely every such surjection lifting $s$ defines such a $\Yt_\varepsilon$. Put $\mathcal{F}:=N_{Y_\varepsilon/\PM_\Dd}\otimes\omega_{Y_\varepsilon/\Dd}$, a locally free sheaf on $Y_\varepsilon$. Tensoring the exact sequence 
$$0 \to \varepsilon \mathcal{O}_Y \to \mathcal{O}_{Y_{\varepsilon}} \to \mathcal{O}_Y \to 0$$
by $\mathcal F$, we have $$0\to\varepsilon\mathcal{F}\to\mathcal{F}\to\NY\otimes K_Y\to 0$$ 
and $\varepsilon\mathcal{F}\cong\NY\otimes K_Y$. By \Cref{lem:pushforward}, $H^0(\NY\otimes K_Y)=k\,s$,  $H^1(\NY\otimes K_Y)=0$. Therefore, $H^0(\mathcal{F})\to H^0(\NY\otimes K_Y)=k\,s$ is surjective with kernel $\varepsilon H^0(\mathcal{F})=\varepsilon k\,s$. So lifts of $s$ exist, are surjective by Nakayama, and form the set $\{(1+\varepsilon\lambda)s_\varepsilon\}$. But every such surjection has the same kernel and therefore we have a unique double structure $\widetilde{Y}_{\varepsilon}$.

(1) With $\eta:=d\Phi(x)$, $I_{\Yt_\varepsilon}\subset I_{Y_\varepsilon}$ means that for $f\in I_{\Yt}$ and any $g$ with $g\equiv\eta(f)$ mod $I_{\Yt}$ one has $g\equiv x(f)$ mod $I_Y$, i.e.\ $\eta(f)\equiv x(f)$ mod $I_Y$ for all $f\in I_{\Yt}$. (2) follows from uniqueness.
\end{proof}

\subsubsection*{The map $\Psi$ from $\mathcal{S}_C$ to the space of canonical ribbons}

Below we define the crucial map $$\Psi: \mathcal{S}_C \to \mathcal{R}_C$$ whose differential will turn out to be the map
\begin{equation*}
\bar\delta^{-1}\circ a\colon\ H^0(\NY(-H))\longrightarrow H^0(N_C\otimes K_Y|_C)\big/\,k\,s_C
\end{equation*}
in \Cref{eq:themap}

\emph{The universal ribbon.} Put $V:=H^0(N_C\otimes L_C)$, so $\RC=\PP(V)$, and let $\pr_2\colon C\times\RC\to\RC$. Since $\pr_{2*}\sHom(N_C^{\vee}\boxtimes\OO_{\RC},L_C\boxtimes\OO_{\RC}(1))=V\otimes\OO_{\RC}(1)$, the tautological line subbundle $\OO_{\RC}(-1)\hookrightarrow V\otimes\OO_{\RC}$, viewed as a nowhere vanishing section of $V\otimes\OO_{\RC}(1)$, is a homomorphism
\[
\tau\colon N_C^{\vee}\boxtimes\OO_{\RC}\longrightarrow L_C\boxtimes\OO_{\RC}(1)
\]
whose restriction to $C\times\{[t]\}$ is a nonzero multiple of $t$; hence $\tau$ is surjective by Lemma~\ref{lem:RC}(1). Let $\Ct_{\RC}\subset H\times\RC$ be the closed subscheme with
\[
I_{\Ct_{\RC}}:=\ker\big(I_{C\times\RC}\to I_{C\times\RC}/I_{C\times\RC}^2=N_C^{\vee}\boxtimes\OO_{\RC}\xrightarrow{\ \tau\ }L_C\boxtimes\OO_{\RC}(1)\big).
\]
From $0\to L_C\boxtimes\OO_{\RC}(1)\to\OO_{\Ct_{\RC}}\to\OO_{C\times\RC}\to0$, $\Ct_{\RC}$ is flat over $\RC$, it is a double structure on $C\times\RC$ with conormal bundle $L_C\boxtimes\OO_{\RC}(1)$, and its fibre over $[t]$ is $\Ct_t$.

\begin{theorem}\label{thm:Psi}
Let $\cY\subset\PM\times\SC$ be the universal scroll and $\pr\colon\cY\to\SC$ the projection; thus $\cY\cap(H\times\SC)=C\times\SC$.
\begin{enumerate}
\item There is a unique closed subscheme $\cYt\subset\PM\times\SC$, flat over $\SC$, with $\cYt_{\mathrm{red}}=\cY$ and all of whose fibres are K3 carpets; its fibre over $[Y']$ is the K3 carpet $\Yt'$ on $Y'$. Moreover $\cN:=\pr_*(N_{\cY}\otimes\omega_{\cY/\SC})$ is a line bundle on $\SC$, $I_{\cY}^2\subset I_{\cYt}$, and $I_{\cY}/I_{\cYt}\cong\omega_{\cY/\SC}\otimes\pr^*\cN^{-1}$.
\item $\cCt:=\cYt\cap(H\times\SC)$ is flat over $\SC$ and is a double structure on $C\times\SC$ in $H\times\SC$ with conormal bundle $I_{C\times\SC}/I_{\cCt}\cong L_C\boxtimes\cN^{-1}$; its fibre over $[Y']$ is the ribbon $\Yt'\cap H$.
\item There is a unique morphism $\Psi\colon\SC\to\RC$ with $(\id_H\times\Psi)^{-1}(\Ct_{\RC})=\cCt$. It satisfies $\Psi^*\OO_{\RC}(-1)\cong\cN$ and, on closed points, $\Psi([Y'])=[\Yt'\cap H]$.
\end{enumerate}
\end{theorem}

\begin{proof}
(1) By Lemma~\ref{lem:SC}(2), $\SC$ is smooth at every $[Y']\in\SC$. Hence \cite[Prop.~3.7]{GP97} applies to the flat family $\pr\colon\cY\to\SC$ of smooth rational normal scrolls in $\PM\times\SC=\PP^g_{\SC}$ and gives the existence and uniqueness of $\cYt$. Its fibre over $[Y']$ is a K3 carpet on $Y'$, hence equals $\Yt'$ \cite[Prop.~1.7]{GP97}. The fact that $\cN$ is a line bundle again follows from \cite[Prop.~1.7]{GP97}. The last two claims follow from the fact that $1\in H^0(\OO_{\SC})=H^0(N_{\cY}\otimes\omega_{\cY/\SC}\otimes\pr^*\cN^{-1})$ is a nowhere vanishing section, and hence gives a surjection $N_{\cY}^{\vee}\to\omega_{\cY/\SC}\otimes\pr^*\cN^{-1}$, and $I_{\cYt}$ is the kernel of $I_{\cY}\to N_{\cY}^{\vee}\to\omega_{\cY/\SC}\otimes\pr^*\cN^{-1}$.

(2) Let $h$ be the equation of $H$. For every $[Y']\in\SC$, $h$ is a nonzerodivisor on $\OO_{Y'}$ and on $\OO_{\Yt'}$ and $Y'$ is smooth, $\Yt'$ is Cohen--Macaulay of pure dimension $2$ with support $Y'\not\subset H$. As $\cY$ and $\cYt$ are flat over $\SC$, the local criterion of flatness \cite[Cor.\ to Thm.~22.5]{Mat89} shows that $h$ is a nonzerodivisor on $\OO_{\cY}$ and on $\OO_{\cYt}$, and that $\OO_{\cCt}=\OO_{\cYt}/h\OO_{\cYt}$ is flat over $\SC$. We have $I_{C\times\SC} = \big(I_{\cY}+(h)\big)$ and $I_{\cCt} = \big(I_{\cYt}+(h)\big)$. Further, $I_{\cY}\cap(h)=hI_{\cY}$. So we have $I_{C\times\SC}^2\subset I_{\cY}^2+(h)\subset I_{\cCt}$ and
\[
I_{C\times\SC}/I_{\cCt}=\big(I_{\cY}+(h)\big)/\big(I_{\cYt}+(h)\big)=I_{\cY}/(I_{\cYt}+hI_{\cY})=(I_{\cY}/I_{\cYt})|_{C\times\SC}\cong\omega_{\cY/\SC}|_{C\times\SC}\otimes\pr^*\cN^{-1},
\]
Finally, $C\times\SC\subset\cY$ is the divisor of $h\in H^0(\OO_{\cY}(1))$, so by adjunction $\omega_{\cY/\SC}|_{C\times\SC}\cong\omega_{C\times\SC/\SC}\otimes\OO_{\cY}(-1)|_{C\times\SC}=(K_C\otimes\OO_C(-1))\boxtimes\OO_{\SC}=L_C\boxtimes\OO_{\SC}$. The fibre of $\cCt$ over $[Y']$ is $\Yt'\cap H$ because base change commutes with fibre products.

(3) By (2), $\cCt$ is defined by the surjection
\[
q\colon N_C^{\vee}\boxtimes\OO_{\SC}=I_{C\times\SC}/I_{C\times\SC}^2\longrightarrow I_{C\times\SC}/I_{\cCt}\cong L_C\boxtimes\cN^{-1},
\]
that is, $I_{\cCt}=\ker\big(I_{C\times\SC}\to N_C^{\vee}\boxtimes\OO_{\SC}\xrightarrow{\ q\ }L_C\boxtimes\cN^{-1}\big)$.
By the projection formula $q\in H^0(\SC,V\otimes\cN^{-1})$, and $q$ vanishes nowhere on $\SC$, its restriction to each $C\times\{[Y']\}$ being a surjection $N_C^{\vee}\to L_C$. Thus $q$ is a line subbundle $\cN\hookrightarrow V\otimes\OO_{\SC}$, i.e.\ a morphism $\Psi\colon\SC\to\PP(V)=\RC$ with $\Psi^*\OO_{\RC}(-1)=\cN$ and $(\id_C\times\Psi)^*\tau=q$. The terms of $0\to I_{\Ct_{\RC}}\to I_{C\times\RC}\to L_C\boxtimes\OO_{\RC}(1)\to0$ are flat over $\RC$, so its pullback along $\id_H\times\Psi$ is the exact sequence $0\to(\id_H\times\Psi)^*I_{\Ct_{\RC}}\to I_{C\times\SC}\xrightarrow{q}L_C\boxtimes\cN^{-1}\to0$; hence $(\id_H\times\Psi)^{-1}(\Ct_{\RC})=\cCt$. Conversely $(\id_H\times\Psi)^{-1}(\Ct_{\RC})$ determines the surjection $N_C^{\vee}\boxtimes\OO_{\SC}\to I_{C\times\SC}/I_{\cCt}$ up to an automorphism of the target, hence the line subbundle $\Psi^*\OO_{\RC}(-1)\subset V\otimes\OO_{\SC}$, hence $\Psi$. On closed points, $\Psi([Y'])$ is the class of the fibre of $\cCt$ over $[Y']$, i.e.\ $[\Yt'\cap H]$ by (2).
\end{proof}

\begin{definition}\label{def:Psi}
Let $x\in H^0(\NY(-H))=T_{[Y]}\SC$, let $Y_\varepsilon\supset C_\Dd$ be the corresponding deformation (Lemma~\ref{lem:SC}(1)), $\Yt_\varepsilon:=\Phi(Y_\varepsilon)$, and
\[
\Ct_\varepsilon:=\Yt_\varepsilon\cap H_\Dd\subset H_\Dd
\]
(scheme-theoretic intersection). Then:
\begin{enumerate}
\item[(i)] $\Ct_\varepsilon$ is flat over $\Dd$ with closed fibre $\Ct$. Indeed $h$ is a nonzerodivisor on $\OO_{\Yt}$, as $\Yt$ is Cohen--Macaulay of pure dimension $2$ with support $Y\not\subset H$; hence $\operatorname{Tor}_1^{\Dd}(\OO_{\Ct_\varepsilon},k)=\ker(h\colon\OO_{\Yt}\to\OO_{\Yt})=0$.
\item[(ii)] $\Ct_\varepsilon$ is a double structure on $C_\Dd$ with conormal $L_C\otimes\OO_\Dd$. Indeed $Y_\varepsilon\cap H_\Dd=C_\Dd$ (it contains $C_\Dd$, and both are flat over $\Dd$ with closed fibre $C$), so $C_\Dd\subset\Ct_\varepsilon$; $I_{C_\Dd}^2\subset I_{\Ct_\varepsilon}$ because $I_{Y_\varepsilon}^2\subset I_{\Yt_\varepsilon}$; and, exactly as in the proof of Lemma~\ref{lem:RC}(1), $I_{C_\Dd/H_\Dd}/I_{\Ct_\varepsilon/H_\Dd}=I_{Y_\varepsilon}/(I_{\Yt_\varepsilon}+hI_{Y_\varepsilon})=L_\varepsilon\otimes\OO_{C_\Dd}\cong L_C\otimes\OO_\Dd$.
\end{enumerate}
Hence, by Lemma~\ref{lem:RC}(2), $\Ct_\varepsilon=\ker(s_C+\varepsilon t)$ for a $t\in H^0(N_C\otimes L_C)$ well defined modulo $k\,s_C$. We define
\[
d\Psi\colon H^0(\NY(-H))\longrightarrow H^0(N_C\otimes L_C)/k\,s_C=T_{[\Ct]}\RC,\qquad x\longmapsto [t].
\]
It is the differential at $[Y]$ of the morphism $\Psi$ of Theorem~\ref{thm:Psi}. Indeed, regard $Y_\varepsilon$ as a morphism $\Dd\to\SC$. By Theorem~\ref{thm:Psi}(1) and flatness, $\cYt\times_{\SC}\Dd$ is flat over $\Dd$, contains $Y_\varepsilon=\cY\times_{\SC}\Dd$, has closed fibre $\Yt$, and satisfies $I_{Y_\varepsilon}^2\subset I_{\cYt\times_{\SC}\Dd}$ with $I_{Y_\varepsilon}/I_{\cYt\times_{\SC}\Dd}=(I_{\cY}/I_{\cYt})|_{Y_\varepsilon}$ invertible; so $\cYt\times_{\SC}\Dd=\Yt_\varepsilon$ by the uniqueness in Lemma~\ref{lem:Phi}. Hence $\Ct_\varepsilon=\cCt\times_{\SC}\Dd$ is the pullback of $\Ct_{\RC}$ along $\Psi\circ Y_\varepsilon\colon\Dd\to\RC$ (Theorem~\ref{thm:Psi}(3)), i.e.\ $\Psi\circ Y_\varepsilon$ is the $\Dd$-point $[s_C+\varepsilon t]$ of $\RC$, which is the tangent vector $[t]\in T_{[\Ct]}\RC$.
\end{definition}

\begin{proposition}\label{prop:a=dPsi}
$a=\bar\delta\circ d\Psi$. Equivalently, the map \eqref{eq:themap} is $d\Psi$.
\end{proposition}

\begin{proof}
Throughout we regard $H^1(\NY(-H+K_Y))$ as a subspace of $H^1(\NYt(-H+K_Y))$ via the injective map $j$ of the bottom row of \eqref{diag:a}; $j$ is induced by the sheaf map $\NY\to\NYt\otimes\OO_Y$, ``restrict a homomorphism from $I_Y$ to $I_{\Yt}$'' (the composite of \eqref{eq:3.3} and \eqref{eq:3.2}), and so is the injection $\iota\colon H^0(\NY(-H))\hookrightarrow H^0(\NYt(-H))$ of \eqref{eq:3.60}, \eqref{eq:3.61}. The maps in play are
\begin{equation}\label{diag:overview}
\begin{tikzcd}[column sep=large, row sep=large]
&& H^0(\NYt(-\Ht))\arrow[d]\\
0\arrow[r] & H^0(\NY(-H))\arrow[r,"\iota"]\arrow[d,dashed,"x\mapsto\eta|_{\Ct}"]\arrow[dr,"a"' near end]\arrow[dd,bend right=60,"d\Psi"'] & H^0(\NYt(-H))\arrow[d,"F"]\\
& H^0(N_C\otimes L_C)\arrow[r,"\delta"]\arrow[d,two heads] & H^1(\NYt(-H+K_Y))\\
& H^0(N_C\otimes L_C)/k\,s_C=T_{[\Ct]}\RC\arrow[ur,"\bar\delta"',"\sim"] &
\end{tikzcd}
\end{equation}
Here $F$ is the connecting map of \eqref{eq:3.65}, i.e.\ of \eqref{eq:3.1} twisted by $-\Ht$, and by definition of $a$ (diagram \eqref{diag:a}) $a=F\circ\iota$. The row through $\delta$ is the cohomology of the restriction sequence $0\to\NY\otimes K_Y(-C)\to\NY\otimes K_Y\to N_C\otimes L_C\to0$ of Lemma~\ref{lem:delta}: $\delta$ is its connecting map, its kernel $k\,s_C$ is the image of $H^0(\NY\otimes K_Y)=k\,s$, and $\bar\delta$ is the induced isomorphism. The dashed arrow is explained below; it is well defined modulo $k\,s_C=\ker\delta$.

Pick $x\in H^0(\NY(-H))$ and set $\xi:=\iota(x)$. \emph{We need to show $\bar\delta(d\Psi(x))=F(\xi)$.}

\emph{Computation of $d\Psi(x)$.} The vector $x$ is a first-order deformation $Y_\varepsilon\subset\PM_\Dd$ of $Y$ containing $C_\Dd$, and $Y_\varepsilon\cap H_\Dd=C_\Dd$ (Lemma~\ref{lem:SC}(1), Definition~\ref{def:Psi}(ii)). By Lemma~\ref{lem:Phi} (the version of \cite[Prop.~3.7]{GP97} over $\Dd$) there is a K3 carpet $\Yt_\varepsilon=\Phi(Y_\varepsilon)\subset\PM_\Dd$ on $Y_\varepsilon$; since $\Yt_\varepsilon\times_\Dd\operatorname{Spec}k=\Yt$, it is an embedded first-order deformation of $\Yt\subset\PM$. Call its class $\eta\in H^0(\NYt)$. Then $\Ct_\varepsilon:=\Yt_\varepsilon\cap H_\Dd$ is a ribbon on $C_\Dd$ with conormal bundle $L_C\otimes\OO_\Dd$, hence a $\Dd$-point of $\RC$ at $[\Ct]$, and by definition (Definition~\ref{def:Psi})
\[
d\Psi(x)=[\Ct_\varepsilon]\in T_{[\Ct]}\RC=H^0(N_C\otimes L_C)/k\,s_C .
\]

\emph{Where $\eta|_{\Ct}$ lives.} Put $N_{\Ct}:=N_{\Ct/H}=\sHom(I_{\Ct},\OO_{\Ct})$. Composing $\eta\colon I_{\Yt}\to\OO_{\Yt}$ with $\OO_{\Yt}\to\OO_{\Ct}$ and using $I_{\Yt}\otimes\OO_H=I_{\Ct/H}$ gives $\eta|_{\Ct}\in H^0(N_{\Ct})$: the class of $\Ct_\varepsilon$ as a first-order deformation of $\Ct$ inside $H$ (reduce the generators $f+\varepsilon\tilde\eta(f)$ of $I_{\Yt_\varepsilon}$ modulo $h$). For the ribbon $\Ct\subset H$, whose conormal bundle is $L_C=I_C/I_{\Ct}$, we have the analogues of \eqref{eq:3.1}--\eqref{eq:3.3} \cite[Lemma~4.1]{GGP08}:
\begin{equation}\label{eq:ribbonseq}
\begin{gathered}
0\to N_{\Ct}\otimes L_C\to N_{\Ct}\to N_{\Ct}\otimes\OO_C\to0,\qquad
0\to\Hc_C\to N_{\Ct}\otimes\OO_C\to L_C^{-2}\to0,\\
0\to L_C^{-1}\xrightarrow{\ s_C^{\vee}\ }N_C\to\Hc_C\to0,\qquad\text{where }\Hc_C:=\sHom(I_{\Ct}/I_C^2,\OO_C).
\end{gathered}
\end{equation}

\emph{Claim.}
\begin{enumerate}
\item[(i)] $\eta|_{\Ct}\in H^0(N_{\Ct}\otimes L_C)\subset H^0(N_{\Ct})$.
\item[(ii)] Restriction of homomorphisms from $I_C$ to $I_{\Ct}$ induces an injection $H^0(N_C\otimes L_C)/k\,s_C\hookrightarrow H^0(N_{\Ct}\otimes L_C)$, and $\eta|_{\Ct}$ lies in its image.
\item[(iii)] Under this injection, $\eta|_{\Ct}=-d\Psi(x)$.
\end{enumerate}

\emph{Proof of the Claim.} (i) The image of $\eta|_{\Ct}$ in $H^0(N_{\Ct}\otimes\OO_C)=\Hom(I_{\Ct},\OO_C)$ is $\eta|_C\colon f\mapsto\eta(f)|_C$. By Lemma~\ref{lem:Phi}(1), $\eta(f)\equiv x(f)$ mod $I_Y$; so $\eta|_C$ is the restriction to $I_{\Ct}$ of the class $f\mapsto x(f)|_C$ of $Y_\varepsilon\cap H_\Dd$ as a deformation of $C$ in $H$, which is $0$ because $Y_\varepsilon\cap H_\Dd=C_\Dd$. Hence $\eta|_C=0$, and by the first sequence of \eqref{eq:ribbonseq}, $\eta|_{\Ct}\in H^0(N_{\Ct}\otimes L_C)=\Hom(I_{\Ct},L_C)$.

(ii) Tensoring the third sequence of \eqref{eq:ribbonseq} with $L_C$ gives $0\to\OO_C\xrightarrow{s_C}N_C\otimes L_C\to\Hc_C\otimes L_C\to0$, whence $H^0(N_C\otimes L_C)/k\,s_C\cong H^0(\Hc_C\otimes L_C)$ as $H^1(\OO_C)=0$; tensoring the second with $L_C$ gives $0\to\Hc_C\otimes L_C\to N_{\Ct}\otimes L_C\to L_C^{-1}\to0$, whence $H^0(\Hc_C\otimes L_C)\hookrightarrow H^0(N_{\Ct}\otimes L_C)$. The composite is $t\mapsto t|_{I_{\Ct}}$. Now $\Ct_\varepsilon$ is a double structure on $C_\Dd$, i.e.\ $I_{C_\Dd}^2\subset I_{\Ct_\varepsilon}$ (Definition~\ref{def:Psi}(ii)); in terms of the class $\eta|_{\Ct}$ this says exactly that $\eta|_{\Ct}\in\Hom(I_{\Ct},L_C)$ vanishes on $I_C^2$, i.e.\ $\eta|_{\Ct}\in\Hom(I_{\Ct}/I_C^2,L_C)=H^0(\Hc_C\otimes L_C)$.

(iii) Write $d\Psi(x)=[t]$, i.e.\ $\Ct_\varepsilon=\ker(s_C+\varepsilon t)$ (Definition~\ref{def:Psi}). For $f\in I_{\Ct}$ the element $f+\varepsilon g$ lies in $I_{\Ct_\varepsilon}$ if and only if $s_C(g)=-t(f)$; so the class of $\Ct_\varepsilon$ is $f\mapsto-t(f)\in L_C=I_C/I_{\Ct}$, that is, $\eta|_{\Ct}=-t|_{I_{\Ct}}$. This proves the Claim.

\emph{The two connecting maps.} Let $\tilde\delta\colon H^0(N_{\Ct}\otimes L_C)\to H^1(\NYt(-H+K_Y))$ be the connecting map of the restriction sequence
\[
0\to\NYt(-H+K_Y)\to\NYt(K_Y)\to\NYt(K_Y)|_C=N_{\Ct}\otimes L_C\to0
\]
of $\NYt(K_Y)=\sHom(I_{\Yt},K_Y)$, $K_Y=I_Y/I_{\Yt}$. Restriction of homomorphisms from $I_Y$ to $I_{\Yt}$ maps the restriction sequence of $\NY\otimes K_Y$ to this one, inducing $j$ on $H^1$ of the first terms and the injection of (ii) on $H^0$ of the third terms; so $\tilde\delta$ restricts to $\delta$ on $H^0(N_C\otimes L_C)/k\,s_C$, and by (iii)
\[
\tilde\delta(\eta|_{\Ct})=-\bar\delta(d\Psi(x)) .
\]
Therefore what we need to show is the following.

\emph{Lemma.} $\tilde\delta(\eta|_{\Ct})=-F(\xi)$.

\emph{Proof of the Lemma.} Multiplication by $h$ followed by restriction to $\Ct$ gives a commutative diagram of sheaves on $\Yt$ with exact rows and columns
\begin{equation}\label{diag:3x3}
\begin{tikzcd}[row sep=small, column sep=small]
& 0\arrow[d] & 0\arrow[d] & 0\arrow[d] & \\
0\arrow[r] & \NYt(-H+K_Y)\arrow[r]\arrow[d,"h"'] & \NYt(-\Ht)\arrow[r]\arrow[d,"h"'] & \NYt(-H)\arrow[r]\arrow[d,"h"] & 0\\
0\arrow[r] & \NYt(K_Y)\arrow[r]\arrow[d] & \NYt\arrow[r]\arrow[d] & \NYt\otimes\OO_Y\arrow[r]\arrow[d] & 0\\
0\arrow[r] & N_{\Ct}\otimes L_C\arrow[r]\arrow[d] & N_{\Ct}\arrow[r]\arrow[d] & N_{\Ct}\otimes\OO_C\arrow[r]\arrow[d] & 0\\
& 0 & 0 & 0 &
\end{tikzcd}
\end{equation}
whose rows are \eqref{eq:3.1} twisted by $-\Ht$, \eqref{eq:3.1} itself, and the first sequence of \eqref{eq:ribbonseq}; the columns are exact because $h$ is a nonzerodivisor on $\OO_{\Yt}$ and on $\OO_Y$. $F$ is the connecting map of the top row and $\tilde\delta$ that of the left column. Lift $\xi$ locally to sections $\xi_i$ of $\NYt(-\Ht)$ over an affine cover $\{U_i\}$; then $F(\xi)=[\xi_i-\xi_j]$. Two observations:
\begin{enumerate}
\item[(a)] View the $\xi_i$ as local sections of $\NYt$ (middle column). Then $\eta|_{U_i}-\xi_i$ restricts to zero on $Y$, since $\eta$ and $\xi_i$ both reduce to $\xi=x|_{I_{\Yt}}$ (Lemma~\ref{lem:Phi}(1)). Hence $\mu_i:=\eta|_{U_i}-\xi_i$ are local sections of $\NYt(K_Y)$.
\item[(b)] The $\xi_i$ vanish on $\Ct$, being multiples of $h$; so $\mu_i|_{\Ct}=\eta|_{\Ct\cap U_i}$, i.e.\ the $\mu_i$ are local lifts of $\eta|_{\Ct}$ to $\NYt(K_Y)$.
\end{enumerate}
Consequently $\tilde\delta(\eta|_{\Ct})=[\mu_i-\mu_j]=[\xi_j-\xi_i]=-F(\xi)$, which proves the Lemma.

By the Lemma, $\bar\delta(d\Psi(x))=-\tilde\delta(\eta|_{\Ct})=F(\xi)=a(x)$.
\end{proof}

\subsection{Proof of Theorem~\ref{thm:alpha}, third part: local sections of \texorpdfstring{$\Psi$}{Psi} and the differential \texorpdfstring{$d\Psi$}{dPsi}}\label{subsec:sections}

\subsubsection*{The bundle $E_{Y'}$}

Put $Q:=K_C\otimes L_C^{-1}=\OO_C(C)=\OO_{\PP^1}(g-1)$. For $Y'\in\SC$ let $\pi'\colon Y'\to C$ be the ruling for which $C$ is a section and $E_{Y'}:=\pi'_*\OO_{Y'}(C)$. Pushing forward we have $0\to\OO_{Y'}\to\OO_{Y'}(C)\to\OO_C(C)\to0$ since $R^1\pi'_*\OO_{Y'}=0$. This gives the extension
\begin{equation}\label{eq:EY}
0\to\OO_C\to E_{Y'}\to Q\to0,
\end{equation}
and $Y'=\PP(E_{Y'})$ with $\OO_{\PP(E_{Y'})}(1)=\OO_{Y'}(C)$ and $C$ the section defined by $E_{Y'}\to Q$. Since $Y'$ is smooth, $E_{Y'}\cong\OO(a_1')\oplus\OO(a_2')$ with $a_i'\ge1$, so \eqref{eq:EY} is not split. Put $V_{Y'}:=E_{Y'}\otimes L_C$, an extension $0\to L_C\to V_{Y'}\to K_C\to0$ with class $[V_{Y'}]\in\Ext^1(K_C,L_C)=\Ext^1(\Omega_C,L_C)$, $[V_{Y'}]\ne0$. We will also denote the extension class of \Cref{eq:EY} by $[E_{Y'}]$.

The two spaces $\RC$ and $\PP(\Ext^1(\Omega_C,L_C))$ are identified by the classical dictionary of Bayer--Eisenbud:

\begin{lemma}\label{lem:beta}
Let $\varepsilon_C\in\Ext^1(\Omega_C,N_C^{\vee})$ be the class of the conormal sequence $0\to N_C^{\vee}\to\Omega_H|_C\to\Omega_C\to0$, and let
\[
\beta\colon H^0(N_C\otimes L_C)=\Hom(N_C^{\vee},L_C)\longrightarrow\Ext^1(\Omega_C,L_C),\qquad t\longmapsto t_*(\varepsilon_C)
\]
(push-out along $t$). Then $\beta(t)$ is the class of the restricted cotangent sequence $0\to L_C\to\Omega_{\Ct_t}|_C\to\Omega_C\to0$ of the double structure $\Ct_t$, and $\beta$ is an isomorphism of $(g-2)$-dimensional vector spaces.
\end{lemma}

\begin{proof}
Tensoring the conormal sequence of $\Ct_t\subset H$ with $\OO_C$ gives $I_{\Ct_t}/I_CI_{\Ct_t}\to\Omega_H|_C\to\Omega_{\Ct_t}|_C\to 0$, and the image of the first map is the image of $\ker t=I_{\Ct_t}/I_C^2\subset N_C^{\vee}$ in $\Omega_H|_C$. Hence $\Omega_{\Ct_t}|_C=\Omega_H|_C/\ker t$, and $0\to N_C^{\vee}/\ker t=L_C\to\Omega_H|_C/\ker t\to\Omega_C\to0$ is the push-out of the conormal sequence along $t$; this is the first claim (cf.\ \cite[Theorem~1.2]{BE95}). For the second, apply $\Hom(-,L_C)$ to the conormal sequence: $\ker\beta$ is the image of $\Hom(\Omega_H|_C,L_C)=H^0(T_H|_C\otimes L_C)$. From the Euler sequence of $H\cong\PP^{g-1}$ restricted to $C$ and tensored with $L_C$,
\[
0\to L_C\to L_C(1)^{\oplus g}\to T_H|_C\otimes L_C\to 0,\qquad L_C(1)=\OO_{\PP^1}(-2),
\]
we get $H^0(T_H|_C\otimes L_C)=\ker\big(H^1(L_C)\to H^1(L_C(1))\otimes H^0(\OO_H(1))^{\vee}\big)$. This map is adjoint to the multiplication pairing $H^1(\OO_{\PP^1}(-g-1))\otimes H^0(\OO_{\PP^1}(g-1))\to H^1(\OO_{\PP^1}(-2))=k$, which is perfect by Serre duality; hence it is injective and $H^0(T_H|_C\otimes L_C)=0$. So $\beta$ is injective, and both spaces have dimension $g-2$.
\end{proof}

\subsubsection*{Local sections of $\Psi$}

\begin{theorem}[Local sections of $\Psi$]\label{thm:Gamma}
Let $\pr\colon C\times\RC\to\RC$ be the projection.
\begin{enumerate}
\item The restricted cotangent sequence of the universal ribbon $\Ct_{\RC}\subset H\times\RC$,
\begin{equation}\label{eq:relcot}
0\to L_C\boxtimes\OO_{\RC}(1)\to\Omega_{\Ct_{\RC}/\RC}|_{C\times\RC}\to K_C\boxtimes\OO_{\RC}\to0,
\end{equation}
is exact, and $\mathcal{V}:=\Omega_{\Ct_{\RC}/\RC}|_{C\times\RC}\otimes(L_C^{-1}\boxtimes\OO_{\RC}(-1))$ sits in an extension
\begin{equation}\label{eq:VR}
0\to\OO_{C\times\RC}\to\mathcal{V}\to Q\boxtimes\OO_{\RC}(-1)\to0
\end{equation}
whose fibre $\mathcal{V}_u$ over $u=[t]\in\RC$ is isomorphic to $\Omega_{\Ct_t}|_C\otimes L_C^{-1}$: a non-split extension of $Q$ by $\OO_C$ with $[\mathcal{V}_u\otimes L_C]\in k^{\times}\beta(t)$.
\item Every point of $\RC$ has an open neighbourhood $U$ with a morphism $\Gamma_U\colon U\to\SC$ such that, for every $u\in U$, $\Gamma_U(u)$ is $\PP(\mathcal{V}_u)$ embedded by $|\OO_{\PP(\mathcal{V}_u)}(1)|$, with $C$ the section defined by $\mathcal{V}_u\to Q$. In particular $[E_{\Gamma_U(u)}]=[\mathcal{V}_u]$ as extensions of $Q$ by $\OO_C$, and $[V_{\Gamma_U(u)}]\in k^{\times}\beta(t)$ for $u=[t]$.
\end{enumerate}
\end{theorem}

\begin{proof}
(1) The sequence \eqref{eq:relcot} is the conormal sequence of $C\times\RC\subset\Ct_{\RC}$ relative to $\RC$: its first term is $I_{C\times\RC}/I_{\Ct_{\RC}}=L_C\boxtimes\OO_{\RC}(1)$ and $\Omega_{C\times\RC/\RC}=K_C\boxtimes\OO_{\RC}$. It is exact on the left because $\Ct_{\RC}$ is locally isomorphic to $(C\times\RC)\times\Dd$ (ribbons are locally split, \cite[\S1]{BE95}; the same local computation applies to $\Ct_{\RC}$, which is defined by the surjection $\tau$). Its restriction to $C\times\{[t]\}$ is the restricted cotangent sequence $0\to L_C\to\Omega_{\Ct_t}|_C\to K_C\to0$ of $\Ct_t$, of class $\beta(t)\ne0$ by Lemma~\ref{lem:beta}. Tensoring with $L_C^{-1}\boxtimes\OO_{\RC}(-1)$ gives \eqref{eq:VR} and the description of its fibres.

(2) The surjection $\mathcal{V}\to Q\boxtimes\OO_{\RC}(-1)$ gives a section $i\colon C\times\RC\to\PP(\mathcal{V})$ of $p\colon\PP(\mathcal{V})\to C\times\RC$ with $i^*\OO_{\PP(\mathcal{V})}(1)=Q\boxtimes\OO_{\RC}(-1)$; since the kernel of this surjection is $\OO_{C\times\RC}$, $i(C\times\RC)$ is the divisor of a section of $\OO_{\PP(\mathcal{V})}(1)$. Let $\mathcal{W}:=\pr_*p_*\OO_{\PP(\mathcal{V})}(1)=\pr_*\mathcal{V}$. For $u\in\RC$, $\mathcal{W}_u=H^0(\mathcal{V}_u)$ with $\mathcal{V}_u\cong\OO(a_1)\oplus\OO(a_2)$, $a_1+a_2=g-1$, $a_i\ge1$ (an extension of $\OO(g-1)$ by $\OO$ with $a_1=0$ is split); so $h^0(\mathcal{V}_u)=g+1=M+1$, $h^1(\mathcal{V}_u)=0$, and $\mathcal{W}$ is a vector bundle of rank $M+1$. Embed $\PP(\mathcal{V})$ in $\PP(\mathcal{W})$ over $\RC$ by $|\OO_{\PP(\mathcal{V})}(1)|$. Pushing forward \eqref{eq:VR},
\[
0\to\OO_{\RC}\to\mathcal{W}\to H^0(Q)\otimes\OO_{\RC}(-1)\to0,
\]
and $i(C\times\RC)$ lies in $\PP\big(H^0(Q)\otimes\OO_{\RC}(-1)\big)=\PP(H^0(Q))\times\RC=H\times\RC$, where it is $C\times\RC$, because $i(C\times\{u\})$ is embedded by $|i^*\OO(1)|=|\OO_C(1)|$. Given $u\in\RC$, let $U\ni u$ be an open set on which $\OO_{\RC}(-1)$ is trivial, e.g.\ the complement of a hyperplane, and choose a splitting of the sequence (it exists, as $\Ext^1(H^0(Q)\otimes\OO_{\RC}(-1),\OO_{\RC})=H^0(Q)^{\vee}\otimes H^1(\OO_{\RC}(1))=0$). This gives $\mathcal{W}|_U\cong(k\oplus H^0(Q))\otimes\OO_U=H^0(\OO_{\PM}(1))\otimes\OO_U$ compatible with the quotients onto $H^0(Q)=H^0(\OO_H(1))$, hence $\PP(\mathcal{W})|_U\cong\PM\times U$ carrying $H\times U$ to $H\times U$ and $C\times U$ to $C\times U$. The composite $\PP(\mathcal{V})|_U\hookrightarrow\PP(\mathcal{W})|_U\cong\PM\times U$ is a flat family of smooth scrolls containing $C\times U$, i.e.\ a morphism $\Gamma_U\colon U\to\SC$. By construction $\Gamma_U(u)=\PP(\mathcal{V}_u)$ embedded by $|\OO_{\PP(\mathcal{V}_u)}(1)|$ with $C=i(C\times\{u\})$; as $\OO_{\PP(\mathcal{V}_u)}(1)=\OO(C)$, $E_{\Gamma_U(u)}=\pi_*\OO_{\Gamma_U(u)}(C)=\mathcal{V}_u$, and \eqref{eq:EY} for $\Gamma_U(u)$ is the fibre of \eqref{eq:VR}. The last assertion follows from (1).
\end{proof}

\subsubsection*{$\GH$-invariance}

\begin{lemma}\label{lem:GH}
\begin{enumerate}
\item $\GH$ acts on $\SC$, and $\Psi\circ\gamma=\Psi$ for every $\gamma\in\GH$.
\item If $Y',Y''\in\SC$ and the classes $[V_{Y'}],[V_{Y''}]\in\Ext^1(K_C,L_C)$ are proportional, then $Y''=\gamma Y'$ for some $\gamma\in\GH$.
\end{enumerate}
\end{lemma}

\begin{proof}
(1) A projectivity $\gamma$ with $\gamma|_H=\id$ maps $Y'\in\SC$ to a smooth scroll $\gamma Y'\supset\gamma C=C$, so $\GH$ acts on $\SC$. Next, $\gamma\Yt'$ is a K3 carpet on $\gamma Y'$, hence is the K3 carpet on $\gamma Y'$ \cite[Prop.~1.7]{GP97}, and $\gamma\Yt'\cap H=\gamma(\Yt'\cap H)=\Yt'\cap H$ as subschemes of $H$ because $\gamma|_H=\id$. Thus $\Psi(\gamma Y')=\Psi(Y')$ on closed points, and $\Psi\circ\gamma=\Psi$ since $\SC$ is reduced (proof of Theorem~\ref{thm:Psi}).

(2) Proportional classes give an isomorphism $E_{Y'}\cong E_{Y''}$ inducing the identity on $Q$ (and a nonzero scalar on $\OO_C$), hence an isomorphism $(\PP(E_{Y'}),C)\cong(\PP(E_{Y''}),C)$ over $C$. So $Y'$ and $Y''$ are two embeddings of the same pair by the complete linear series $|\OO(C)|$ which agree on $C$; they differ by a $\gamma\in PGL_{M+1}(k)$ with $\gamma|_C=\id$, hence $\gamma|_H=\id$ since $C$ spans $H$, i.e.\ $\gamma\in\GH$.
\end{proof}

\subsubsection*{$\Gamma_U$ is a local section of $\Psi$}

\begin{proposition}\label{prop:BEsection}
Let $\Gamma_U\colon U\to\SC$ be as in Theorem~\ref{thm:Gamma}(2).
\begin{enumerate}
\item $\Psi\circ\Gamma_U=\id_U$.
\item For every $Y'\in\SC$, $\Psi(Y')=[t]$ for the $t\in H^0(N_C\otimes L_C)$ with $\beta(t)\in k^{\times}[V_{Y'}]$, i.e, the hyperplane section $\Yt'\cap H$ of the carpet is the ribbon whose restricted cotangent extension is that of $[V_{Y'}]=[E_{Y'}\otimes L_C]$.
\item $\Gamma_U(\Psi(Y'))\in\GH\cdot Y'$ for every $Y'\in\Psi^{-1}(U)$.
\end{enumerate}
\end{proposition}

\begin{proof}
(1) Let $u=[t]\in U$ and $\Ct:=\Ct_t\subset H$. It is a non-split ribbon on $C$ with conormal bundle $L_C$ (Lemma~\ref{lem:RC}, Lemma~\ref{lem:beta}), and it is canonically embedded in $H=\PP^{g-1}$: it is nondegenerate of degree $2g-2$ and arithmetic genus $g$ in $\PP^{g-1}$, so by Riemann--Roch $\omega_{\Ct}\otimes\OO_{\Ct}(-1)$ is a line bundle of degree $0$ with a nonzero section, hence trivial (its restriction to $C$ is trivial, and $H^0(L_C)=0$). We apply Theorem~\ref{thm:extension}(3) to $\Ct\subset H$ (so $g_C=0$ in Set-up~\ref{setup:ribbon}). By Theorem~\ref{thm:Gamma}(1) the restricted cotangent sequence of $\Ct$ is the fibre of \eqref{eq:relcot} at $u$, so $V=\Omega_{\Ct}|_C=\mathcal V_u\otimes L_C$; hence the ruled surface of Set-up~\ref{setup:ribbon} is $\PP(V)=\PP(\mathcal V_u)$, with $\OO_{\PP(V)}(C)=\OO_{\PP(V)}(1)\otimes p^*L_C^{-1}=\OO_{\PP(\mathcal V_u)}(1)$ (Lemma~\ref{lem:OSC}(1)) and $C$ the section defined by $\mathcal V_u\to Q$. By Theorem~\ref{thm:Gamma}(2), $\Gamma_U(u)\subset\PM$ is $\PP(\mathcal V_u)$ embedded by the complete linear series of $\OO_{\PP(\mathcal V_u)}(1)$, with $H$ the hyperplane such that $\Gamma_U(u)\cap H=C$ is the given curve (Lemma~\ref{lem:SC}); and $\Yt'$ is the K3 carpet on $\Gamma_U(u)$. Theorem~\ref{thm:extension}(3) gives $\Yt'\cap H=\Ct$, hence $\Psi(\Gamma_U(u))=[\Yt'\cap H]=[\Ct]=u$.

(2), (3) Let $Y'\in\SC$. Since $[E_{Y'}]$ is non-split, $[V_{Y'}]\ne0$, and since $\beta$ is an isomorphism (Lemma~\ref{lem:beta}) there is a $t\ne0$ with $\beta(t)=[V_{Y'}]$. Put $u:=[t]$ and choose $\Gamma_{U'}$ with $u\in U'$ (Theorem~\ref{thm:Gamma}(2)). Then $[V_{\Gamma_{U'}(u)}]\in k^{\times}\beta(t)=k^{\times}[V_{Y'}]$, so $Y'\in\GH\cdot\Gamma_{U'}(u)$ by Lemma~\ref{lem:GH}(2), and by Lemma~\ref{lem:GH}(1) and part (1),
\[
\Psi(Y')=\Psi(\Gamma_{U'}(u))=u=[t] .
\]
This is (2). If moreover $\Psi(Y')\in U$, then $u=\Psi(Y')\in U$, and $[V_{\Gamma_U(u)}]$ also lies in $k^{\times}\beta(t)$, so $\Gamma_U(u)\in\GH\cdot\Gamma_{U'}(u)=\GH\cdot Y'$ by Lemma~\ref{lem:GH}(2). This is (3).
\end{proof}

\subsubsection*{The differential of $\Psi$}

\begin{proposition}\label{prop:dPsi}
Let $d\Psi\colon H^0(\NY(-H))=T_{[Y]}\SC\to T_{[\Ct]}\RC=H^0(N_C\otimes L_C)/k\,s_C$ be the differential of $\Psi$ at $[Y]$ (Definition~\ref{def:Psi}).
\begin{enumerate}
\item $d\Psi$ is surjective.
\item $\ker d\Psi=\Im c=T_{[Y]}(\GH\cdot Y)$.
\end{enumerate}
\end{proposition}

\begin{proof}
(1) Choose an open $U\ni\Psi(Y)$ and $\Gamma_U$ as in Theorem~\ref{thm:Gamma}(2). By Proposition~\ref{prop:BEsection}(1), $\Psi\circ\Gamma_U=\id_U$; differentiating at $\Psi(Y)$,
\[
(d\Psi)_{\Gamma_U(\Psi(Y))}\circ {(d\Gamma_U)}_{\Psi(Y)}=\id\colon T_{\Psi(Y)}\RC\to T_{\Psi(Y)}\RC,
\]
so $d\Psi$ is surjective at $\Gamma_U(\Psi(Y))=\gamma Y$, $\gamma\in\GH$ (Proposition~\ref{prop:BEsection}(3)). By Lemma~\ref{lem:GH}(1), $d\Psi_Y=d\Psi_{\gamma Y}\circ d\gamma_Y$, and $d\gamma_Y\colon T_Y\SC\to T_{\gamma Y}\SC$ is an isomorphism since $\gamma$ is an automorphism of $\SC$; hence $d\Psi_Y$ is surjective.

(2) By Lemma~\ref{lem:GH}(1), the two morphisms $\GH\times\SC\to\RC$, $(\gamma,Y')\mapsto\Psi(\gamma Y')$ and $(\gamma,Y')\mapsto\Psi(Y')$, agree on closed points, and hence agree on $\GH\times\SC$ since it is reduced. Differentiating at $(1,Y)$ along $\GH$ gives $d\Psi\circ c=0$, i.e.\ $\Im c\subset\ker d\Psi$. By (1) and Lemma~\ref{lem:SC},
\[
\dim\ker d\Psi=h^0(\NY(-H))-\dim\RC=(M+1)+(2b-e-2)-(2b-e-2)=M+1=\dim\Im c ,
\]
so $\ker d\Psi=\Im c$.
\end{proof}

\subsection{End of the proof of Theorem~\ref{thm:alpha}}

\begin{proposition}\label{prop:a-surj}
The map $\bar\delta^{-1}\circ a=d\Psi\colon H^0(\NY(-H))\to H^0(N_C\otimes K_Y|_C)/k\,s_C$ is surjective. Consequently $a$ and $F$ are surjective, and $\ker a=\Im c$.
\end{proposition}

\begin{proof}
By Proposition~\ref{prop:dPsi}, $d\Psi$ is surjective with kernel $\Im c$. By Proposition~\ref{prop:a=dPsi}, $a=\bar\delta\circ d\Psi$ is surjective with kernel $\Im c$. Surjectivity of $F$ follows from diagram \eqref{diag:a} and the sequence \eqref{eq:3.61}, as explained in the proof of Theorem~\ref{thm:alpha}.
\end{proof}

\begin{proof}[Proof of Theorem~\ref{thm:alpha}]
By \eqref{eq:alpha-F}, $\alpha(\Yt)=h^0(-H-2K_Y)+\dim\coker F$, and $F$ is surjective by Proposition~\ref{prop:a-surj}. For the last assertion note that $-H-2K_Y=3C_0+(2e+4-b)f$ and $h^0(\OO_Y(3C_0+kf))=\sum_{i=0}^{3}h^0(\OO_{\PP^1}(k-ie))$, which vanishes if and only if $k<0$, i.e.\ $b\ge 2e+5$.
\end{proof}

\begin{remark}[Deformation-theoretic meaning]
The equality $\ker a=\Im c$ obtained in Proposition~\ref{prop:a-surj} says: a first-order deformation $Y_\varepsilon$ of the scroll inside $\PM$ which contains $C_\Dd$ is contained in a first-order deformation $\Yt_\varepsilon$ of the carpet with $\Yt_\varepsilon\cap H_\Dd=\Ct_\Dd$ if and only if $Y_\varepsilon=\gamma_\varepsilon Y$ for a projectivity $\gamma_\varepsilon\in\GH(\Dd)$ fixing $H$ pointwise; and then $\Yt_\varepsilon=\gamma_\varepsilon\Yt$. Equivalently, by Proposition~\ref{prop:BEsection}(2): the hyperplane section of the K3 carpet on a scroll $Y'\supset C$ is the ribbon whose restricted cotangent bundle is $V_{Y'}=E_{Y'}\otimes L_C$ ($\cong E_{Y'}^{\vee}\otimes K_C$), and this ribbon determines $Y'$ up to $\GH$. The scalar $\yt_h$ is the ``horizontal component'' of the carpet class of \cite[Proposition~1.7]{GP97}.
\end{remark}

\section{Prime K3 surfaces of low genus and Fano threefolds of Picard rank one}\label{sec:fano}

This section is the case $r=1$ of \cite[\S4]{BM25}; the bounds on $\alpha$ are now supplied by Theorem~\ref{thm:alpha} (and by Proposition~\ref{prop:g6} for genus $6$), and the vanishing of $h^0(N(-2))$ by Proposition~\ref{prop:N2}.

\subsection{Set-up}\label{subsec:fano-setup}

For a smooth K3 surface $X\subset\PP^n$ we write $N_X:=N_{X/\PP^n}$ and, as in \S\ref{subsec:setup}, $\alpha(X):=h^0(N_X(-1))-(n+1)$. Recall \cite{Zak91,Lvo92} that if $\alpha(X)\le0$ then $X$ is not extendable: every $V\subset\PP^{n+1}$ having $X$ as a hyperplane section is a cone over $X$.

Fix $g\ge3$. Let $\mathcal H_g$ be the locus, in the Hilbert scheme of $\PP^g$, of smooth K3 surfaces $X\subset\PP^g$ embedded by the complete linear series of a very ample line bundle $H$ with $H^2=2g-2$ and $H$ primitive in $\operatorname{Pic}X$; we call them \emph{prime K3 surfaces of genus $g$}. Since the moduli space of primitively polarized K3 surfaces of genus $g$ is irreducible \cite[Ch.~6]{Huy16}, $\mathcal H_g$ is irreducible; it is smooth of dimension $18+(g+1)^2$, because for $X\in\mathcal H_g$ the Euler sequence gives $h^0(T_{\PP^g}|_X)=(g+1)^2-1$ and $H^1(T_{\PP^g}|_X)=H^2(\OO_X)=k$, the map $H^1(T_X)\to H^2(\OO_X)$, cup product with $c_1(H)$, is surjective, and $H^2(T_X)=0$, so that $h^1(N_X)=0$ and $h^0(N_X)=(g+1)^2-1+19$. A point of $\mathcal H_g$ is called \emph{general} if it lies in a suitable dense open subset.

Let $e\in\{0,1\}$ and $b\ge e+2$, and let $\Yt\subset\PP^g$ be the K3 carpet on $Y=F_e$ embedded by $|H|$, $H=C_0+bf$, so that $g=2b-e+1$ (this is the set-up of \S\ref{subsec:setup}, except that we allow $b=e+2$). By \cite[Theorems~3.5 and 4.4]{GP97}, $[\Yt]$ lies in the closure of $\mathcal H_g$. Since $\Yt$ and the smooth K3 surfaces are local complete intersections, the normal sheaf of a flat family in $\PP^g$ degenerating a general $X\in\mathcal H_g$ to $\Yt$ is locally free and commutes with base change; by upper semicontinuity, for $X$ general in $\mathcal H_g$ and every $k\ge0$,
\begin{equation}\label{eq:semicont}
h^0(N_X(-k))\le h^0(\NYt(-k\Ht)),\qquad\text{in particular}\qquad \alpha(X)\le\alpha(\Yt).
\end{equation}

\subsection{The invariant \texorpdfstring{$\alpha$}{alpha} for \texorpdfstring{$6\le g\le 12$}{6 <= g <= 12}}

\begin{theorem}\label{thm:alpha-lowgenus}
Let $X\subset\PP^g$ be a general prime K3 surface of genus $g$.
\begin{enumerate}
\item If $g=12$, then $\alpha(X)\le1$.
\item If $g=10$, then $\alpha(X)\le3$.
\item If $g=9$, then $\alpha(X)\le4$.
\item If $g=8$, then $\alpha(X)\le6$.
\item If $g=7$, then $\alpha(X)\le8$.
\item If $g=6$, then $\alpha(X)\le10$.
\item If $g=11$ or $g\ge13$, then $\alpha(X)\le0$; consequently $X$ is not extendable.
\end{enumerate}
\end{theorem}

\begin{proof}
We can degenerate such a K3 surface to a K3 carpet $\Yt\subset\PP^g$, extending $Y\subset\PP^g$, where $Y=F_e$ is embedded by the complete linear series of $C_0+bf$ with $e=0$ or $e=1$ depending on whether $g$ is odd or even respectively, and $g=2b-e+1$. By \eqref{eq:semicont} it is enough to find an upper bound for $\alpha(\Yt)$. In cases (1)--(5) the pair $(e,b)$ is $(1,6)$, $(1,5)$, $(0,4)$, $(1,4)$, $(0,3)$ respectively, and in case (7) it is $(0,b)$ with $b\ge5$ or $(1,b)$ with $b\ge7$. In all these cases $b\ge e+3$, so by Theorem~\ref{thm:alpha}
\[
\alpha(\Yt)=h^0\big(\OO_Y(3C_0+(2e+4-b)f)\big)=\sum_{i=0}^{3}h^0\big(\OO_{\PP^1}(2e+4-b-ie)\big),
\]
which equals $1,3,4,6,8$ in cases (1)--(5) and $0$ in case (7). In case (6), $(e,b)=(1,3)$ and $\alpha(\Yt)=10$ by Proposition~\ref{prop:g6} below. The last assertion in (7) is the theorem of Zak and L'vovsky.
\end{proof}

\begin{proposition}[The case $g=6$]\label{prop:g6}
Let $Y=F_1\subset\PP^6$ be embedded by $|C_0+3f|$ and let $\Yt$ be the K3 carpet on $Y$. Then $\alpha(\Yt)=h^0(\OO_Y(-H-2K_Y))=h^0(\OO_Y(3C_0+3f))=10$.
\end{proposition}

\begin{proof}
Here $b=e+2$, $M=6$, and $\OO_Y(-H-K_Y)=\OO_Y(C_0)$, with $h^0(\OO_Y(C_0))=1$ and $h^1(\OO_Y(C_0))=h^2(\OO_Y(C_0))=0$. We go through the proof of Theorem~\ref{thm:alpha}, where the assumption $b\ge e+3$ was used only to have $H^0(-H-K_Y)=0$ in \eqref{eq:3.41} and \eqref{eq:3.52}; everything else, in particular \eqref{eq:3.47}, \eqref{eq:3.48}, \eqref{eq:3.56}--\eqref{eq:3.59} and \eqref{eq:3.61}, holds for $b\ge e+1$. From \eqref{eq:3.41} and \eqref{eq:3.48} we get an exact sequence
\begin{multline*}
0\to H^0(\NYt(-H+K_Y))\to H^0(\OO_Y(C_0))=k\xrightarrow{\ \delta_1\ }H^1(\NY(-H+K_Y))=k^3\\
\xrightarrow{\ j\ }H^1(\NYt(-H+K_Y))\to0
\end{multline*}
and $H^2(\NYt(-H+K_Y))=0$; thus $h^0(\NYt(-H+K_Y))=1-\operatorname{rank}\delta_1$, $h^1(\NYt(-H+K_Y))=3-\operatorname{rank}\delta_1$, and $j$ is surjective. From \eqref{eq:3.52} and \eqref{eq:3.58}, $h^0(\Hc(-H))=h^0(\NY(-H))-1=(M+1)+h^1(T_Y(-H))-1=7+3-1=9$, and from \eqref{eq:3.61}, $h^0(\NYt(-H))=9+h^0(\OO_Y(3C_0+3f))=19$ and $H^1(\NYt(-H))=H^1(\OO_Y(3C_0+3f))=0$. Sequence \eqref{eq:3.40} therefore gives
\[
h^0(\NYt(-\Ht))=h^0(\NYt(-H+K_Y))+h^0(\NYt(-H))-\operatorname{rank}F=20-\operatorname{rank}\delta_1-\operatorname{rank}F,
\]
where $F\colon H^0(\NYt(-H))\to H^1(\NYt(-H+K_Y))$ is the connecting map, and it remains to show that $F$ is surjective: then $\operatorname{rank}F=3-\operatorname{rank}\delta_1$ and $\alpha(\Yt)=17-7=10$.

The proof of Proposition~\ref{prop:a=dPsi} shows, without using $b\ge e+3$, that $F(\iota(x))=j(\delta(t))$ for every $x\in H^0(\NY(-H))$, where $\iota\colon H^0(\NY(-H))\to H^0(\NYt(-H))$ is induced by restriction of homomorphisms from $I_Y$ to $I_{\Yt}$, $d\Psi(x)=[t]$, and $\delta$ is the map of Lemma~\ref{lem:delta}. Now $\delta$ is surjective (Lemma~\ref{lem:delta}), $d\Psi$ is surjective (Proposition~\ref{prop:dPsi}; its proof, and those of the results of \S\S\ref{subsec:geometric}--\ref{subsec:sections} on which it relies, only use $b\ge e+1$), and $j$ is surjective; hence the image of $F$ is all of $H^1(\NYt(-H+K_Y))$.
\end{proof}

\subsection{Sections of \texorpdfstring{$N(-2)$}{N(-2)}}

\begin{proposition}\label{prop:N2}
Let $Y=F_e\subset\PM$ be embedded by $|H|$, $H=C_0+bf$ with $b\ge e+2$, and let $\Yt$ be the K3 carpet on $Y$. Then
\[
h^0(\NYt(-2\Ht))\le h^0(\OO_Y(-2H-2K_Y))=h^0\big(\OO_Y(2C_0+(2e+4-2b)f)\big).
\]
In particular $h^0(\NYt(-2\Ht))=0$ if $b\ge e+3$, and $h^0(\NYt(-2\Ht))\le1$ if $(e,b)=(1,3)$.
\end{proposition}

\begin{proof}
Tensoring \eqref{eq:3.1} by $\OO_{\Yt}(-2\Ht)$ gives $h^0(\NYt(-2\Ht))\le h^0(\NYt(-2H+K_Y))+h^0(\NYt(-2H))$. We use throughout that $H^0(\OO_Y(-H+K_Y))=0$, $H^0(\OO_Y(-2H-K_Y))=0$ (as $-2H-K_Y=(e+2-2b)f$ and $2b>e+2$), and, by Kodaira vanishing, $H^1(\OO_Y(-2H))=H^1(\OO_Y(-2H+K_Y))=0$ and and $H^i(\OO_Y(-H))=0$ for all $i$ (see the derivation of \eqref{eq:3.54}).

\emph{Step 1: $H^0(\NYt(-2H+K_Y))=0$.} By \eqref{eq:3.2} tensored with $-2H+K_Y$, $H^0(\NYt(-2H+K_Y))=H^0(\Hc(-2H+K_Y))$, and by \eqref{eq:3.3} tensored with $-2H+K_Y$ the latter is a quotient of $H^0(\NY(-2H+K_Y))$. By \eqref{eq:3.11} tensored with $-2H+K_Y$, $H^0(\TP(-2H+K_Y))=0$, so by \eqref{eq:3.10} $H^0(\NY(-2H+K_Y))$ injects into $H^1(T_Y(-2H+K_Y))$. Finally $H^1(T_Y(-2H+K_Y))=0$: in
\[
0\to T_{Y/\PP^1}(-2H+K_Y)\to T_Y(-2H+K_Y)\to p^*T_{\PP^1}(-2H+K_Y)\to0
\]
we have $T_{Y/\PP^1}(-2H+K_Y)=\OO_Y(-2C_0-(2b+2)f)$ and $p^*T_{\PP^1}(-2H+K_Y)=\OO_Y(-4C_0-(2b+e)f)$; both have $p_*=0$, while by relative duality $R^1p_*\OO_Y(-2C_0)=\OO_{\PP^1}(e)$ and $R^1p_*\OO_Y(-4C_0)=\OO_{\PP^1}(e)\oplus\OO_{\PP^1}(2e)\oplus\OO_{\PP^1}(3e)$. Hence
\begin{gather*}
H^1(T_{Y/\PP^1}(-2H+K_Y))=H^0(\OO_{\PP^1}(e-2b-2))=0,\\
H^1(p^*T_{\PP^1}(-2H+K_Y))=H^0\big(\OO_{\PP^1}(-2b)\oplus\OO_{\PP^1}(e-2b)\oplus\OO_{\PP^1}(2e-2b)\big)=0 .
\end{gather*}

\emph{Step 2: $H^0(\Hc(-2H))=0$.} Granting this, \eqref{eq:3.2} tensored with $-2H$ gives $h^0(\NYt(-2H))\le h^0(\OO_Y(-2H-2K_Y))$, which together with Step~1 proves the inequality. By \eqref{eq:3.3} tensored with $-2H$ we have an exact sequence
\[
0\to H^0(\NY(-2H))\to H^0(\Hc(-2H))\to H^1(\OO_Y(-2H-K_Y))\xrightarrow{\ s^{\vee}\ }H^1(\NY(-2H)),
\]
so it suffices to show that $H^0(\NY(-2H))=0$ and that $s^{\vee}$ is injective.

(a) By \eqref{eq:3.11} tensored with $-2H$, $H^0(\TP(-2H))=0$ and $H^1(\TP(-2H))\cong H^2(\OO_Y(-2H))$. Hence by \eqref{eq:3.10}, $H^0(\NY(-2H))=\ker\big(H^1(T_Y(-2H))\to H^1(\TP(-2H))\big)$. In $0\to T_{Y/\PP^1}(-2H)\to T_Y(-2H)\to p^*T_{\PP^1}(-2H)\to0$ the sheaf $p^*T_{\PP^1}(-2H)=\OO_Y(-2C_0+(2-2b)f)$ has $H^0=H^1=0$ ($p_*=0$ and $R^1p_*=\OO_{\PP^1}(e+2-2b)$), so $H^1(T_{Y/\PP^1}(-2H))\to H^1(T_Y(-2H))$ is an isomorphism. Put $E:=p_*\OO_Y(H)=\OO_{\PP^1}(b)\oplus\OO_{\PP^1}(b-e)$, so that $Y=\PP(E)$ with $\OO_{\PP(E)}(1)=\OO_Y(H)$ and $H^0(\OO_Y(H))=H^0(E)$. The relative Euler sequence $0\to\OO_Y\to p^*E^{\vee}\otimes\OO_Y(H)\to T_{Y/\PP^1}\to0$ maps to the Euler sequence \eqref{eq:3.11} via the dual of the evaluation map $H^0(E)\otimes\OO_{\PP^1}\to E$, the induced map $T_{Y/\PP^1}\to\TP$ being the natural inclusion. Twisting by $-2H$, the connecting map $H^1(T_{Y/\PP^1}(-2H))\to H^2(\OO_Y(-2H))$ of the relative Euler sequence therefore factors through $H^1(\TP(-2H))\cong H^2(\OO_Y(-2H))$; and it is injective, because $H^1(p^*E^{\vee}\otimes\OO_Y(-H))=0$ ($p^*E^{\vee}\otimes\OO_Y(-H)=p^*E^{\vee}(-b)\otimes\OO_Y(-C_0)$ and $p_*\OO_Y(-C_0)=R^1p_*\OO_Y(-C_0)=0$). Hence $H^1(T_Y(-2H))\to H^1(\TP(-2H))$ is injective and $H^0(\NY(-2H))=0$. Moreover $h^1(T_Y(-2H))=h^1(\OO_Y((e-2b)f))=2b-e-1=h^0(\OO_Y((2b-e-2)f))=h^2(\OO_Y(-2H))$, so this map is an isomorphism, and \eqref{eq:3.10} shows that the connecting map $\partial\colon H^1(\NY(-2H))\to H^2(T_Y(-2H))$ is injective.

(b) It suffices to show that $\partial\circ s^{\vee}\colon H^1(\OO_Y(-2H-K_Y))\to H^2(T_Y(-2H))$ is injective. The pull-back of the normal sequence \eqref{eq:3.10} along $s^{\vee}\colon\OO_Y(-K_Y)\to\NY$ is an extension of $\OO_Y(-K_Y)$ by $T_Y$ whose class is $\partial(s)=\yt\in\Ext^1(\OO_Y(-K_Y),T_Y)=H^1(T_Y\otimes K_Y)$ (Lemma~\ref{lem:pushforward}); so $\partial\circ s^{\vee}$ is the cup product with $\yt$. Write $\OO_Y(-2H-K_Y)=p^*L$ with $L:=\OO_{\PP^1}(e+2-2b)$. By Lemma~\ref{lem:pushforward} (and its proof) and the projection formula, $p_*(T_Y\otimes K_Y\otimes p^*L)=L(-2)$ and $R^1p_*(T_Y\otimes K_Y\otimes p^*L)=L$, so the Leray spectral sequence identifies $H^2(T_Y(-2H))=H^2(T_Y\otimes K_Y\otimes p^*L)$ with $H^1(\PP^1,L)=H^1(\OO_Y(-2H-K_Y))$, and by its multiplicativity the cup product with $\yt$ becomes multiplication by the image of $\yt$ in $H^0(R^1p_*(T_Y\otimes K_Y))=H^0(\OO_{\PP^1})=k$, which is nonzero by Lemma~\ref{lem:pushforward}. Hence $\partial\circ s^{\vee}$ is an isomorphism, $s^{\vee}$ is injective, and $H^0(\Hc(-2H))=0$.

Finally $h^0(\OO_Y(2C_0+(2e+4-2b)f))=\sum_{i=0}^{2}h^0(\OO_{\PP^1}(2e+4-2b-ie))$, which is $0$ for $b\ge e+3$ and $1$ for $(e,b)=(1,3)$.
\end{proof}

\subsection{Classification of Fano threefolds of Picard rank one}

\begin{definition}\label{def:Vg}
For $g\ge3$ let $\mathcal V_g$ denote the locus, in the Hilbert scheme of $\PP^{g+1}$, of smooth Fano threefolds $V\subset\PP^{g+1}$ of Picard rank one, index one and genus $g$, i.e.\ $\operatorname{Pic}V=\mathbb Z(-K_V)$ and $(-K_V)^3=2g-2$, embedded by the complete linear series $|-K_V|$; it is open in the Hilbert scheme, and $\overline{\mathcal V_g}$ denotes its closure. For $X\subset\PP^g$ we denote by $C(X)\subset\PP^{g+1}$ the cone over $X$.
\end{definition}

\begin{corollary}\label{cor:cone}
Let $X\subset\PP^g$ be a general prime K3 surface of genus $g\ge6$. Then the tangent space to the Hilbert scheme of $\PP^{g+1}$ at $[C(X)]$ has dimension
\[
\dim T_{[C(X)]}=h^0(N_X)+h^0(N_X(-1))+h^0(N_X(-2))=18+(g+1)^2+(g+1)+\alpha(X)+h^0(N_X(-2)),
\]
where $h^0(N_X(-2))=0$ if $g\ge7$ and $h^0(N_X(-2))\le1$ if $g=6$.
\end{corollary}

\begin{proof}
Since $X$ is projectively normal and $H^1(\OO_X(k))=0$ for all $k$, one has $H^0(N_{C(X)/\PP^{g+1}})=\bigoplus_{k\ge0}H^0(N_X(-k))$ (see \cite[\S2]{CLM93}). We have $h^0(N_X)=18+(g+1)^2$ (\S\ref{subsec:fano-setup}) and $h^0(N_X(-1))=(g+1)+\alpha(X)$ by definition. Since $\operatorname{Pic}X=\mathbb ZH$ and $g\ge5$, the homogeneous ideal of $X$ is generated by quadrics \cite{SD74}; hence $N_X=\sHom(I_X/I_X^2,\OO_X)$ is a subsheaf of a direct sum of copies of $\OO_X(2)$, and $H^0(N_X(-k))=0$ for $k\ge3$. Finally $h^0(N_X(-2))\le h^0(\NYt(-2\Ht))$ by \eqref{eq:semicont}, and the right-hand side is $0$ for $g\ge7$ (then $b\ge e+3$) and at most $1$ for $g=6$ ($(e,b)=(1,3)$) by Proposition~\ref{prop:N2}.
\end{proof}

\begin{lemma}\label{lem:hypsection}
Let $W$ be an irreducible component of $\overline{\mathcal V_g}$ and $V$ a general element of $W$. Then a general hyperplane section of $V$ corresponds to a general K3 surface in $\mathcal H_g$.
\end{lemma}

\begin{proof}
The proof follows from standard deformation theory arguments (see for example \cite[Lemma~3.7]{CLM98}).
\end{proof}

\begin{theorem}\label{thm:fano}
Let $g\ge6$.
\begin{enumerate}
\item $\mathcal V_g=\emptyset$ if $g=11$ or $g\ge13$.
\item For $6\le g\le10$ and $g=12$, $\mathcal V_g$ is irreducible, and the families of Fano and Iskovskikh \cite{Isk77,Isk78} (see \cite{Muk88} and \cite[Table~3.1]{CLM98}) form an open dense subset of $\mathcal V_g$.
\end{enumerate}
\end{theorem}

\begin{proof}
\emph{Proof of (1).} If $\mathcal V_g\neq\emptyset$, let $V$ be a general element of an irreducible component of $\overline{\mathcal V_g}$. By Lemma~\ref{lem:hypsection} a general hyperplane section of $V$ is a general prime K3 surface $X$ of genus $g$, and $V$ is not a cone; this contradicts Theorem~\ref{thm:alpha-lowgenus}(7).

\emph{Proof of (2).} First one notes that $\mathcal H_g$ is irreducible. Now the list of examples of Fano--Iskovskikh--Mori--Mukai (see for example \cite[Table~3.1]{CLM98}) shows that in each of the cases mentioned in part (2), there exist K3 surfaces in $\mathcal H_g$ which are extendable to smooth Fano threefolds. For a smooth Fano threefold $V\in\mathcal V_g$ one has $h^1(N_{V/\PP^{g+1}})=0$ (from \eqref{eq:3.10} for $V$: $H^1(T_{\PP^{g+1}}|_V)=0$ by the Euler sequence, as $H^1(\OO_V(1))=H^2(\OO_V)=0$, and $H^2(T_V)=0$ by Kodaira--Nakano vanishing), so $V$ is a smooth point of the Hilbert scheme and lies on a unique irreducible component, of dimension $h^0(N_{V/\PP^{g+1}})$. Hence the examples of Fano--Iskovskikh--Mori--Mukai, which form an irreducible family, lie on a unique irreducible component $W_0$ of $\overline{\mathcal V_g}$, and by Lemma~\ref{lem:hypsection} $\mathcal H_g$ is dominated by $W_0$. We end the proof by showing that $W_0$ is the unique irreducible component of $\overline{\mathcal V_g}$ and that the known examples form an open dense subset of $\mathcal V_g$. To see this, we observe first, as in \cite{CLM93} and \cite{CLM98}, that a projectively Cohen--Macaulay scheme degenerates along a flat family to the cone over its hyperplane section. Hence, by Lemma~\ref{lem:hypsection} applied to each irreducible component $W$ of $\overline{\mathcal V_g}$, every $W$ contains the cone $C(X)$ over a general K3 surface $X\subset\PP^g$ in $\mathcal H_g$, and we may take the same $X$ for all components, $\mathcal H_g$ being irreducible; so it is enough to show that for a general $X\in\mathcal H_g$ the cone $C(X)\subset\PP^{g+1}$ is a smooth point of the Hilbert scheme. For the stated values of $g$, Corollary~\ref{cor:cone} along with Theorem~\ref{thm:alpha-lowgenus} gives an upper bound for $\dim T_{[C(X)]}$, which is listed as follows:
\begin{center}
\begin{tabular}{c|cccccc}
$g$ & $6$ & $7$ & $8$ & $9$ & $10$ & $12$\\ \hline
$\dim T_{[C(X)]}\le$ & $85$ & $98$ & $114$ & $132$ & $153$ & $201$
\end{tabular}
\end{center}
Now using the descriptions of the examples of Fano--Iskovskikh--Mori--Mukai, one can calculate the dimension $h^0(N_{V/\PP^{g+1}})$ (see for example \cite[Table~2, p.~666]{CLM93} and \cite[Table~3.1]{CLM98}) to check that in each case $\dim W_0=h^0(N_{V/\PP^{g+1}})$ is exactly the upper bound just obtained. Since $C(X)\in W_0$, we get $\dim W_0\le\dim T_{[C(X)]}\le\dim W_0$; hence $C(X)$ is a smooth point of the Hilbert scheme, and $W_0$ is the only component through it.
\end{proof}

\section{Extensions of ribbons to \texorpdfstring{$K$}{K}-trivial carpets}\label{sec:extensions}

This section is \S7 of \cite{DM25}, reproduced with the additions needed to complete the proof of Theorem~7.5(1) of that paper (Theorem~\ref{thm:extension}(1) below) and with the embedded versions of that statement (Theorem~\ref{thm:extension}(2),(3)). The notation is that of \S\ref{sec:carpets} with the curve $C$ allowed to have any genus $g_C$: $\Ct$ is a ribbon on $C$ of arithmetic genus $g$ with conormal bundle $L_C$, $Y$ is a ruled surface over $C$ and $\Yt$ a carpet on $Y$; for $g_C=0$ one recovers the objects of \S\ref{sec:carpets}. For a vector bundle $V$ on $C$ we write $\PP(V)=\operatorname{Proj}\operatorname{Sym}V$, so that $p_*\OO_{\PP(V)}(1)=V$ and a surjection $V\to A$ onto a line bundle defines a section $i$ of $p\colon\PP(V)\to C$ with $i^*\OO_{\PP(V)}(1)=A$. A \emph{ribbon} on $C$ with conormal bundle $L_C$ is a double structure $\Ct$ on $C$ with $I_{C/\Ct}^2=0$ and $I_{C/\Ct}\cong L_C$; a \emph{$K$-trivial carpet} on $Y$ is a double structure $\Yt$ on $Y$ with $I_{Y/\Yt}^2=0$ and $I_{Y/\Yt}\cong K_Y$ (for $g_C=0$ these are the K3 carpets of \S\ref{sec:carpets}). By \cite[Theorem~1.2]{BE95} (whose proof applies verbatim to carpets, see \cite[Lemma~3.1]{DM25}), ribbons on $C$ with conormal bundle $L_C$, together with an identification $I_{C/\Ct}=L_C$, correspond bijectively to $\Ext^1(\Omega_C,L_C)$ via the restricted cotangent sequence
\[
0\to L_C\to\Omega_{\Ct}|_C\to\Omega_C\to0 ,
\]
two ribbons being isomorphic as ribbons on $C$ (i.e.\ by an isomorphism inducing the identity on $C$) if and only if their classes are proportional; the class $0$ is the split ribbon. Likewise $K$-trivial carpets on $Y$ with an identification $I_{Y/\Yt}=K_Y$ correspond to $\Ext^1(\Omega_Y,K_Y)$ via $0\to K_Y\to\Omega_{\Yt}|_Y\to\Omega_Y\to0$; we write $\Yt_{\yt}$ for the carpet of class $\yt$, so that $\Yt_{\lambda\yt}\cong\Yt_{\yt}$ over $Y$ for $\lambda\in k^{\times}$ and $\Yt_0$ is the split carpet.

\begin{setup}\label{setup:ribbon}
Let $\Ct$ be a ribbon of arithmetic genus $g$ and conormal bundle $L_C$ over a smooth curve $C$ of genus $g_C$. Then $\Ct$ is defined by an extension
\[
0\to L_C\to V\to K_C\to0
\]
where $V$ is a rank two vector bundle on $C$. Further $V=\Omega_{\Ct}|_C$. Now the data of the extension canonically gives
\begin{enumerate}
\item a ruled surface $Y=\PP(V)\to C$ and
\item an embedding $C\hookrightarrow Y=\PP(V)$ as a section of the ruled surface $Y$ defined by the surjection $V\to K_C$.
\end{enumerate}
Note that $\deg L_C=2g_C-1-g$, since $\chi(\OO_{\Ct})=\chi(\OO_C)+\chi(L_C)$.
\end{setup}

\begin{definition}\label{def:invariant}
Every ruled surface on a fixed curve $C$ can be uniquely expressed as $Y=\PP(V_n)$ where $V_n$ is a normalized vector bundle, i.e.\ $h^0(V_n)\neq0$ and $h^0(V_n\otimes A^{-1})=0$ for every line bundle $A$ of positive degree. The invariant of a ruled surface is defined as $e=-\deg(V_n)$. Let $C_0$ denote the numerical class of $\OO_{\PP(V_n)}(1)$ and $f$ denote the numerical class of any fiber.
\end{definition}

\begin{lemma}\label{lem:invariant}
In Set-up~\ref{setup:ribbon}, assume that $g>g_C+1$. Then
\begin{enumerate}
\item the invariant of the ruled surface $Y$ is
\[
e=-\deg(L_C)+(2g_C-2)-2b(\Ct)=g-2b(\Ct)-1
\]
where $b(\Ct)$ is the blow-up index of $\Ct$ (see \cite{DM24}).
\item If $\Ct$ is a ribbon on $C$ with arithmetic genus $g$, such that $g>\max\{2g_C-1,g_C+1\}$ and maximum blow-up index, then
\begin{enumerate}
\item if both $g_C$ and $g$ are even, $e=-(g_C-1)$
\item if $g_C$ is even and $g$ is odd, $e=-g_C$
\item if $g_C$ is odd and $g$ is even, $e=-g_C$
\item if both $g_C$ and $g$ are odd, $e=-(g_C-1)$.
\end{enumerate}
Hence $Y$ is general with respect to the invariant $e$.
\item If $\Ct$ on $C$ is a general ribbon with arithmetic genus $g$, such that $g>\max\{4g_C-3,g_C+1\}$, the ruled surface $Y$ is general in its space of deformations.
\end{enumerate}
\end{lemma}

\begin{proof}
\emph{Proof of (1).} The normalized vector bundle $V_n$ is given by $V_n=V\otimes A^{-1}$ where $A$ is the line sub-bundle of maximum degree. By \cite[Proposition~4.4]{DM24}, we have that $\deg(A)=2g_C-2-b(\Ct)$. Hence
\begin{align*}
-\deg(V_n)&=-\deg(V\otimes A^{-1})\\
&=-\deg(L_C\otimes A^{-1})-\deg(K_C\otimes A^{-1})\\
&=g-(2g_C-1)-(2g_C-2)+2(2g_C-2-b(\Ct))\\
&=g-2b(\Ct)-1 .
\end{align*}

\emph{Proof of (2).} This follows by plugging in the maximum value of the blow-up index from \cite[Corollary~4.3]{DM24}.

\emph{Proof of (3).} Let $0\to L_C\to V\to K_C\to0$ be an extension defining a general ribbon. We show that for a one parameter family of deformations $\mathcal V$ on $C\times T\to T$ of $V_0=V$, the general fiber $V_t$ admits a surjection $V_t\to K_C$. Note that the surjection $V\to K_C\to0$ gives a global section of $H^0(V^*\otimes K_C)$. Now $h^1(V^*\otimes K_C)=h^0(V)$ by Serre duality. For a general ribbon, we can assume $V$ to be semistable (see \cite[Corollary~4.5]{DM24}). From the defining extension, we have $\mu(V)=\frac12\deg(K_C\otimes L_C)=\frac12(2g_C-2-g+(2g_C-1))=\frac12(4g_C-3-g)<0$ by the condition on the arithmetic genus. Hence $h^1(V^*\otimes K_C)=h^0(V)=0$. This means the section of $V^*\otimes K_C$ deforms along the one-parameter family to give a section of $V_t^*\otimes K_C$ for a fiber over a general $t$. This section $V_t\to K_C$ is surjective by Nakayama's lemma.
\end{proof}

\begin{lemma}\label{lem:OSC}
In Set-up~\ref{setup:ribbon},
\begin{enumerate}
\item The line bundle $\OO_Y(C)=\OO_{\PP(V)}(1)\otimes p^*L_C^{-1}$. The numerical class of this line bundle is given by $C_0+(g-b(\Ct)-1)f=C_0+\frac12(g+e-1)f$.
\item Assume that $V$ is semistable (which holds if $b(\Ct)\ge\frac12(g+3)-2$, or equivalently $e\le0$, by \cite[Corollary~4.5]{DM24}) and that $g>4g_C+1$. Then the multiplication map
\[
H^0(\OO_Y(C))\otimes H^0(\OO_Y(C))\to H^0(\OO_Y(2C))
\]
is surjective.
\end{enumerate}
\end{lemma}

\begin{proof}
\emph{Proof of (1).} The section $C$ is the zero locus of the composite $p^*L_C\to p^*V\to\OO_{\PP(V)}(1)$, which vanishes exactly at the points where the quotient of $V$ is $V\to K_C$; hence $\OO_Y(C)=\OO_{\PP(V)}(1)\otimes p^*L_C^{-1}$. We have that $V_n=V\otimes A^{-1}$ where $A$ is the maximum degree line subbundle of $V$. Hence $\OO_{\PP(V)}(1)=\OO_{\PP(V_n)}(1)\otimes p^*A$, and $\OO_Y(C)=\OO_{\PP(V)}(1)\otimes p^*L_C^{-1}=\OO_{\PP(V_n)}(1)\otimes p^*(A\otimes L_C^{-1})$. Now $\deg(A\otimes L_C^{-1})=\deg(A)-\deg(L_C)=2g_C-2-b(\Ct)+g-(2g_C-1)=g-b(\Ct)-1=\frac12(g+e-1)$.

\emph{Proof of (2).} We need to show that the map $H^0(\OO_{\PP(V)}(1)\otimes p^*L_C^{-1})\otimes H^0(\OO_{\PP(V)}(1)\otimes p^*L_C^{-1})\to H^0(\OO_{\PP(V)}(2)\otimes p^*L_C^{-2})$ is surjective. Pushing forward we see that the map is actually a composition of the two following maps on the curve $C$:
\[
H^0(V\otimes L_C^{-1})\otimes H^0(V\otimes L_C^{-1})\to H^0(V^{\otimes2}\otimes L_C^{-2})\to H^0(\operatorname{Sym}^2V\otimes L_C^{-2}) .
\]
Considering that $\operatorname{Sym}^2V\otimes L_C^{-2}$ is a direct summand of $V^{\otimes2}\otimes L_C^{-2}$, it is enough to show the surjectivity of the multiplication map of global sections of vector bundles on the curve $C$. The vector bundle $V=\Omega_{\Ct}|_C$ is semistable by assumption. Then by \cite[Theorem~1]{But94}, the surjection holds if $\mu(V\otimes L_C^{-1})>2g_C$. This happens if $\deg(K_C\otimes L_C^{-1})>4g_C$ or in other words $-\deg(L_C)>2g_C+2$. This is equivalent to $g>4g_C+1$.
\end{proof}

\begin{theorem}\label{thm:extension}
In Set-up~\ref{setup:ribbon}, assume that $\Ct$ is non-split. For $\yt\in\Ext^1(\Omega_Y,K_Y)$ let $\Yt_{\yt}$ be the $K$-trivial carpet on $Y$ of class $\yt$.
\begin{enumerate}
\item There is a unique $\yt_0\in\Ext^1(\Omega_Y,K_Y)$ such that the embedding $i\colon C\hookrightarrow Y$ extends to a closed embedding $\Ct\hookrightarrow\Yt_{\yt_0}$ inducing the identity $K_Y|_C=L_C$ on conormal bundles.
\item Suppose $\OO_Y(C)$ is very ample, let $Y\hookrightarrow\PP^M$ be the embedding by the complete linear series of $\OO_Y(C)$, $M+1=h^0(\OO_Y(C))$, let $H\subset\PP^M$ be the hyperplane with $Y\cap H=C$, and let $\Yt\hookrightarrow\PP^M$ be the unique $K$-trivial carpet extending the embedding $Y\hookrightarrow\PP^M$ \cite[Proposition~3.2]{DM25}. Then $\Yt\cap H$ is a ribbon on $C$ with conormal bundle $L_C$, and $\Yt\cap H\cong\Ct$ as ribbons on $C$.
\item Let $g_C=0$ and let $\Ct\subset\PP^{M-1}$ be a canonically embedded ribbon, $M=g$, so that $C\subset\PP^{M-1}$ is a rational normal curve. Embed $Y$ in $\PM$ by the complete linear series of $\OO_Y(C)$ and identify $\PP^{M-1}$ with a hyperplane $H\subset\PM$ so that $Y\cap H=C$. If $\Yt\subset\PM$ is the K3 carpet on $Y$, then
\[
\Yt\cap H=\Ct
\]
in $\PP^{M-1}=H$.
\end{enumerate}
\end{theorem}

\begin{proof}
\emph{Proof of (1).} Let $C$ be a smooth curve and $L_C$ a line bundle on $C$. Fix a ribbon $\Ct$ on $C$ with conormal bundle $L_C$. Then $\Ct$ defines an element $[\Ct]\in\Ext^1(\Omega_C,L_C)$. Set $V=\Omega_{\Ct}|_C$. Then we have the exact sequence
\[
0\to L_C\to V\to\Omega_C\to0 .
\]
Let $Y=\PP V$. The exact sequence above gives an embedding
\[
i\colon C\hookrightarrow Y
\]
which is a section of the projection $p\colon Y\to C$. We have
\[
K_Y=\OO_{\PP(V)}(-2)\otimes p^*(K_C\otimes\det V)
\]
and hence
\[
K_Y|_C=\Omega_C^{-2}\otimes(\Omega_C\otimes L_C)\otimes K_C=L_C .
\]
A $K$-trivial carpet $\Yt$ on $Y$ is given by an element $\yt\in\Ext^1(\Omega_Y,K_Y)$. To show the existence of a $K$-trivial carpet as claimed in part (1), we want to choose $\Yt=\Yt_{\yt}$ so that the embedding
\[
C\xhookrightarrow{\ i\ }Y
\]
extends to
\[
\Ct\xhookrightarrow{\ \tilde i\ }\Yt
\]
inducing the identity on conormal bundles. This is equivalent to the equality $\bar i[\Ct]=r(\yt)$ in the following diagram:
\begin{equation}\label{diag:ext}
\begin{tikzcd}
& \Ext^1(\Omega_Y,K_Y)\arrow[d,"r"]\\
\Ext^1(\Omega_C,L_C)\arrow[r,"\bar i"] & \Ext^1(\Omega_Y|_C,L_C)
\end{tikzcd}
\end{equation}
where $r$ is the map induced by the restriction of the short exact sequence $0\to K_Y\to\Omega_{\Yt}|_Y\to\Omega_Y\to0$ to $C$ and $\bar i\colon\Ext^1(\Omega_C,L_C)\to\Ext^1(\Omega_Y|_C,L_C)$ is the map induced by the embedding $i\colon C\hookrightarrow Y$.

Now the projection $p\colon Y\to C$ induces a map $p^*\Omega_C\to\Omega_Y$ which when restricted to the section $i\colon C\hookrightarrow Y$ induces a map $\Omega_C\to\Omega_Y|_C$. This further induces the map $\bar p\colon\Ext^1(\Omega_Y|_C,L_C)\to\Ext^1(\Omega_C,L_C)$ which is a left inverse to the map $\bar i$: $\bar p\circ\bar i=\id$. Thus the conormal sequence $0\to N_{C/Y}^{\vee}\to\Omega_Y|_C\to\Omega_C\to0$ is split by $\Omega_C\to\Omega_Y|_C$, and
\begin{equation}\label{eq:extsplit}
\Ext^1(\Omega_Y|_C,L_C)=\Ext^1(N_{C/Y}^{\vee},L_C)\oplus\Ext^1(\Omega_C,L_C),
\end{equation}
where $\bar i$ is the inclusion of the second summand, $\bar p$ the projection onto it, and the projection onto the first summand is the map
\[
r_N\colon\Ext^1(\Omega_Y|_C,L_C)\to\Ext^1(N_{C/Y}^{\vee},L_C)
\]
induced by $N_{C/Y}^{\vee}\hookrightarrow\Omega_Y|_C$. So it suffices to prove that there is a unique $\yt_0\in\Ext^1(\Omega_Y,K_Y)$ satisfying the two conditions
\begin{equation}\label{eq:twoconditions}
\text{(a)}\ \ [\Ct]=\bar p\,r(\yt_0),\qquad\qquad\text{(b)}\ \ r_N\,r(\yt_0)=0 .
\end{equation}

\emph{Condition (a).} We have to describe the $\yt$ for which $[\Ct]$ is the image of $\yt$ under the map
\[
\Ext^1(\Omega_Y,K_Y)\xrightarrow{\ r\ }\Ext^1(\Omega_Y|_C,L_C)\xrightarrow{\ \bar p\ }\Ext^1(\Omega_C,L_C)
\]
which also factors as
\[
\Ext^1(\Omega_Y,K_Y)\to\Ext^1(p^*\Omega_C,K_Y)\to\Ext^1(\Omega_C,L_C).
\]
We first show that the map $\Ext^1(\Omega_Y,K_Y)\to\Ext^1(p^*\Omega_C,K_Y)$ is surjective. Considering the exact sequence
\[
0\to p^*\Omega_C\to\Omega_Y\to\Omega_{Y/C}\to0 ,
\]
it is enough to show that $\Ext^2(\Omega_{Y/C},K_Y)=0$. By Serre duality,
\[
\Ext^2(\Omega_{Y/C},K_Y)=\Hom(\OO_Y,\Omega_{Y/C})^{\vee}=0
\]
since $\Omega_{Y/C}$ is $\OO_{\PP(V)}(-2)\otimes p^*(\det V)$. First note that since $p_*K_Y=0$, we have
\[
\Ext^1(p^*\Omega_C,K_Y)=H^1(K_Y\otimes p^*T_C)=H^0(R^1p_*K_Y\otimes T_C)=H^0(\OO_C),
\]
because
\[
R^1p_*K_Y=R^1p_*(K_{Y/C}\otimes p^*K_C)=K_C .
\]
So $\Ext^1(p^*\Omega_C,K_Y)=k$. We claim that the image of $1\in\Ext^1(p^*\Omega_C,K_Y)=H^1(\Omega_{Y/C})$ (note that $K_Y\otimes p^*T_C=\Omega_{Y/C}$) in $\Ext^1(\Omega_C,L_C)=H^1(T_C\otimes L_C)$ is $[\Ct]$. This follows from the following general claim.

\begin{claim}\label{claim:euler}
Let $0\to L\to V\to A\to0$ be an extension where $L$ and $A$ are line bundles on $C$. Consider the section $i\colon C\hookrightarrow\PP(V)$ defined by the extension. This induces a map $\Omega_{\PP(V)/C}\to\Omega_{\PP(V)/C}|_C$. Then the image of $1\in H^1(\Omega_{\PP(V)/C})=k$ inside $H^1(\Omega_{\PP(V)/C}|_C)=H^1(L\otimes A^{-1})$ is the class of the extension.
\end{claim}

\emph{Proof of the claim.} $1\in H^1(\Omega_{\PP V/C})$ is represented by the Euler sequence
\[
0\to\Omega_{\PP V/C}\to p^*V\otimes\OO_{\PP V}(-1)\to\OO_{\PP V}\to0 .
\]
Restricted to $C$ via $i$, it is
\[
0\to L\otimes A^{-1}\to V\otimes A^{-1}\to\OO_C\to0 ,
\]
which is the class of the extension twisted by $A^{-1}$. This proves the claim.

For $\yt\in\Ext^1(\Omega_Y,K_Y)$ let $\yt_h\in\Ext^1(p^*\Omega_C,K_Y)=k$ be its image. By the claim,
\begin{equation}\label{eq:pbar-r}
\bar p\,r(\yt)=\yt_h\,[\Ct] ,
\end{equation}
and since $\Ct$ is non-split, $[\Ct]\ne0$; so condition (a) holds for $\yt$ if and only if $\yt_h=1$. The kernel of the map $\Ext^1(\Omega_Y,K_Y)\to\Ext^1(p^*\Omega_C,K_Y)$ is
\[
\Ext^1(\Omega_{Y/C},K_Y)=H^1(T_{Y/C}\otimes K_Y)=H^1(p^*K_C)=H^1(K_C)=k .
\]
Hence the set of $\yt$ satisfying condition (a) is the affine line
\[
B:=\{\yt\in\Ext^1(\Omega_Y,K_Y):\ \yt_h=1\},
\]
the inverse image of $1\in\Ext^1(p^*\Omega_C,K_Y)$.

\emph{Condition (b).} We have a diagram as follows:
\begin{equation}\label{diag:B}
\begin{tikzcd}
\Ext^1(\Omega_{Y/C},K_Y)\arrow[r]\arrow[d] & \Ext^1(\Omega_Y,K_Y)\arrow[r]\arrow[d,"r"] & \Ext^1(p^*\Omega_C,K_Y)=k\arrow[d]\\
\Ext^1(\Omega_{Y/C}|_C,L_C)\arrow[r,hook] & \Ext^1(\Omega_Y|_C,L_C)\arrow[r,"\bar p"] & \Ext^1(\Omega_C,L_C)
\end{tikzcd}
\end{equation}
whose vertical maps are restriction to $C$. Let $B\subset\Ext^1(\Omega_Y,K_Y)$ be the affine line above, and let $K$ be the kernel of $r_N\colon\Ext^1(\Omega_Y|_C,L_C)\to\Ext^1(N_{C/Y}^{\vee},L_C)=k$. We show that $r(B)\cap K$ consists of a single point, i.e.\ that $r_N\circ r$ vanishes at exactly one point of $B$.

Two distinct elements of $B$ differ by a nonzero element of $\Ext^1(\Omega_{Y/C},K_Y)$. This is because
\[
\Hom(p^*\Omega_C,K_Y)=H^0(K_Y\otimes p^*\Omega_C^{*})=0 .
\]
Now note that the conormal sequence
\[
0\to N_{C/Y}^{\vee}\to\Omega_Y|_C\to\Omega_C\to0
\]
has a splitting given by
\[
0\to p^*\Omega_C|_C=\Omega_C\to\Omega_Y|_C\to\Omega_{Y/C}|_C\to0 ,
\]
the restriction to $C$ of $0\to p^*\Omega_C\to\Omega_Y\to\Omega_{Y/C}\to0$. Therefore $\Omega_{Y/C}|_C\cong N_{C/Y}^{\vee}$, and $N_{C/Y}^{\vee}\hookrightarrow\Omega_Y|_C\to\Omega_{Y/C}|_C$ is the identity $N_{C/Y}^{\vee}\to N_{C/Y}^{\vee}$. Hence, for $w\in\Ext^1(\Omega_{Y/C},K_Y)$, $r_N\,r(w)$ is the image of $w$ under the left vertical map of \eqref{diag:B},
\[
\Ext^1(\Omega_{Y/C},K_Y)\to\Ext^1(\Omega_{Y/C}|_C,L_C)=\Ext^1(N_{C/Y}^{\vee},L_C),
\]
and it is enough to show that this map is injective. Tensoring
\[
0\to\OO_Y(-C)\to\OO_Y\to\OO_C\to0
\]
with $\Omega_{Y/C}^{*}\otimes K_Y$ we get
\[
0\to\Omega_{Y/C}^{*}\otimes K_Y\otimes\OO_Y(-C)\to\Omega_{Y/C}^{*}\otimes K_Y\to\Omega_{Y/C}^{*}|_C\otimes L_C\to0 .
\]
Since $\Omega_{Y/C}=K_Y\otimes p^*K_C^{*}$, we have $\Omega_{Y/C}^{*}=K_Y^{*}\otimes p^*K_C$, and hence, by Lemma~\ref{lem:OSC}(1),
\[
\Omega_{Y/C}^{*}\otimes K_Y\otimes\OO_Y(-C)=p^*K_C\otimes\OO_Y(-C)=\OO_{\PP V}(-1)\otimes p^*(K_C\otimes L_C),
\]
whose $H^1$ vanishes because $R^jp_*\OO_{\PP V}(-1)=0$ for all $j$. Taking cohomology of the last exact sequence gives the injectivity of $H^1(\Omega_{Y/C}^{*}\otimes K_Y)\to H^1(\Omega_{Y/C}^{*}|_C\otimes L_C)$, i.e.\ of the left vertical map of \eqref{diag:B}.

Therefore $r_N\circ r$ is injective on $B$. Being an injective affine map from the line $B$ to the line $\Ext^1(N_{C/Y}^{\vee},L_C)=k$, it is bijective, so exactly one point $\yt_0$ of $B$ satisfies condition (b). This proves (1).

\emph{Proof of (2).} Note that since by Lemma~\ref{lem:OSC}, part (1), $\OO_Y(C)=\OO_{\PP(V)}(1)\otimes p^*L_C^{-1}$, we have the existence of a unique $K$-trivial carpet $\Yt\hookrightarrow\PP^M$ extending the embedding $Y\hookrightarrow\PP^M$, and it is non-split \cite[Proposition~3.2]{DM25}; let $\yt\ne0$ be its class. Let $h$ be a linear form defining $H$, so that $h|_Y$ is the section of $\OO_Y(C)$ with zero scheme $C$. Since $h$ is a nonzerodivisor on $\OO_Y$ and on $K_Y$, it is a nonzerodivisor on $\OO_{\Yt}$, and $K_Y\cap h\OO_{\Yt}=hK_Y$. Hence the kernel of $\OO_{\Yt\cap H}=\OO_{\Yt}/h\OO_{\Yt}\to\OO_Y/h\OO_Y=\OO_C$ is $K_Y/hK_Y=K_Y\otimes\OO_C=L_C$, and it has square zero. Thus $\Ct':=\Yt\cap H$ is a ribbon on $C$ with conormal bundle $L_C$, $\Ct'\cap Y=C$, and the map $K_Y|_C\to I_{C/\Ct'}=L_C$ induced by $\Ct'\subset\Yt$ is the identity. As in the proof of (1), $r(\yt)=\bar i[\Ct']$; hence, by \eqref{eq:extsplit} and \eqref{eq:pbar-r}, $r_N\,r(\yt)=0$ and $[\Ct']=\yt_h[\Ct]$. By the proof of condition (b), $r_N\circ r$ is injective on $\Ext^1(\Omega_{Y/C},K_Y)=\ker(\yt\mapsto\yt_h)$; as $\yt\ne0$ and $r_N\,r(\yt)=0$, we get $\yt_h\ne0$. Therefore $[\Ct']$ and $[\Ct]$ are proportional and nonzero, and $\Ct'\cong\Ct$ as ribbons on $C$.

\emph{Proof of (3).} Since $\Ct$ is non-split, $V\otimes L_C^{-1}\cong\OO_{\PP^1}(a)\oplus\OO_{\PP^1}(a')$ with $a,a'\ge1$ (an extension $0\to\OO_C\to V\otimes L_C^{-1}\to K_C\otimes L_C^{-1}\to0$ with $a=0$ is split), so $\OO_Y(C)=\OO_{\PP(V\otimes L_C^{-1})}(1)$ is very ample; and $h^0(\OO_Y(C))=M+1$, since $H^1(\OO_Y)=0$ and $N_{C/Y}=K_C\otimes L_C^{-1}=\OO_{\PP^1}(g-1)$.

Let $\Yt_{\yt_0}$ and $\Ct\subset\Yt_{\yt_0}$ be as in (1). Then $\Ct$ is a Cartier divisor on $\Yt_{\yt_0}$ with $\OO_{\Yt_{\yt_0}}(\Ct)|_Y=\OO_Y(C)$, and from
\[
0\to K_Y(C)\to\OO_{\Yt_{\yt_0}}(\Ct)\to\OO_Y(C)\to0
\]
and $H^i(K_Y(C))=H^{2-i}(\OO_Y(-C))^{\vee}=0$ for all $i$ (Lemma~\ref{lem:OSC}(1): $\OO_Y(-C)=\OO_{\PP V}(-1)\otimes p^*L_C$) we get $H^0(\OO_{\Yt_{\yt_0}}(\Ct))\cong H^0(\OO_Y(C))$. Let
\[
\varphi\colon\Yt_{\yt_0}\to\PP^M_1:=\PP\big(H^0(\OO_{\Yt_{\yt_0}}(\Ct))^{\vee}\big)
\]
be the morphism defined by the complete linear series of $\OO_{\Yt_{\yt_0}}(\Ct)$. Its restriction to $Y$ is the embedding by $|\OO_Y(C)|$, and $\varphi^{-1}(H_1)=\Ct$ for the hyperplane $H_1$ defined by the section of $\OO_{\Yt_{\yt_0}}(\Ct)$ vanishing on $\Ct$.

We claim that $\varphi$ is a closed embedding. It is finite (proper and injective on points), so it is enough to show that $\OO_{\PP^M_1}\to\varphi_*\OO_{\Yt_{\yt_0}}$ is surjective. Its image is a subsheaf of algebras mapping onto $\OO_Y=\varphi_*\OO_{\Yt_{\yt_0}}/K_Y$, so it is everything if and only if it contains $K_Y$, i.e.\ if and only if the map $I_Y\to K_Y$ induced by $\OO_{\PP^M_1}\to\varphi_*\OO_{\Yt_{\yt_0}}$ is surjective. Since $K_Y^2=0$ this map factors through $N_Y^{\vee}=I_Y/I_Y^2$, so it is a section of $N_Y\otimes K_Y$; and $H^0(N_Y\otimes K_Y)=k\,s$ with $s$ nowhere vanishing \cite[Proposition~3.2]{DM25}. Thus the map is $\lambda s$ for some $\lambda\in k$, and it is surjective if $\lambda\ne0$. If $\lambda=0$, then $I_Y$ maps to $0$, so $\OO_{\PP^M_1}\to\varphi_*\OO_{\Yt_{\yt_0}}$ factors through $\OO_Y$, i.e.\ $\varphi$ factors through a retraction of $Y\subset\Yt_{\yt_0}$; a retraction splits $0\to K_Y\to\OO_{\Yt_{\yt_0}}\to\OO_Y\to0$ as $\OO_Y$-algebras, so $\Yt_{\yt_0}$ would be the split carpet, contradicting $(\yt_0)_h=1$. Hence $\varphi$ is a closed embedding; $\varphi(Y)\subset\PP^M_1$ is embedded by $|\OO_Y(C)|$, $\varphi(\Yt_{\yt_0})$ is a K3 carpet on it, and $\varphi(\Yt_{\yt_0})\cap H_1=\varphi(\Ct)$.

It is clear that $\varphi(\Ct)\subset H_1=\PP^{g-1}$ is canonically embedded.

Now the given $\Ct\subset\PP^{M-1}$ and $\varphi(\Ct)\subset H_1$ are two embeddings of $\Ct$ by the complete linear series of $\omega_{\Ct}$, so there is a projectivity $T\colon H_1\to\PP^{M-1}$ with $T\circ\varphi|_{\Ct}$ the given inclusion; extend $T$ to a projectivity $\hat T\colon\PP^M_1\to\PM$ with $\hat T(H_1)=H$. Then $\hat T\circ\varphi|_Y$ and the given embedding $Y\subset\PM$ are two embeddings of $Y$ by the complete linear series $|\OO_Y(C)|$ which agree on $C$, so they differ by a projectivity $\gamma$ of $\PM$ with $\gamma|_C=\id$, hence $\gamma|_H=\id$ as $C$ spans $H$. Finally $\gamma\hat T\varphi(\Yt_{\yt_0})$ is a K3 carpet on $Y\subset\PM$, hence is $\Yt$ \cite[Proposition~1.7]{GP97}, and
\[
\Yt\cap H=\gamma\hat T\big(\varphi(\Yt_{\yt_0})\cap H_1\big)=\gamma\hat T\varphi(\Ct)=\gamma(\Ct)=\Ct .
\]
\end{proof}

\begin{corollary}\label{cor:nonsmoothable}
In Set-up~\ref{setup:ribbon}, let $g_C\ge2$, let $\Ct$ be non-split and let $\Yt=\Yt_{\yt_0}$ be the carpet of Theorem~\ref{thm:extension}(1). Then
\begin{enumerate}
\item $\Yt$ is not smoothable.
\item If $b(\Ct)\ge\frac12(g+3)-2$ or equivalently $e\le0$, every infinitesimal deformation of $\Yt$ is locally trivial.
\end{enumerate}
\end{corollary}

\begin{proof}
Part (1) follows immediately from \cite[Theorem~5.1(1)]{DM25}. To show part (2), we only need to show that under the conditions on the blow-up index, or equivalently the invariant $e$, the vector bundle $V$ of the defining extension
\[
0\to L_C\to V\to K_C\to0
\]
is semistable, and apply \cite[Theorem~5.1(2)]{DM25}. This follows from \cite[Corollary~4.5]{DM24}.
\end{proof}

\begin{remark}\label{rem:BE}
Theorem~\ref{thm:extension}(3) says that a canonically embedded non-split ribbon $\Ct\subset\PP^{g-1}$ of arithmetic genus $g$ on the rational normal curve $C$ is the hyperplane section $\Yt\cap H$ of the K3 carpet on the rational normal scroll $Y=\PP(\Omega_{\Ct}|_C\otimes L_C^{-1})\subset\PP^{g}$, embedded by $|\OO_{\PP(\Omega_{\Ct}|_C\otimes L_C^{-1})}(1)|=|\OO_Y(C)|$, $H=\PP^{g-1}$ being the hyperplane with $Y\cap H=C$. This is a theorem of Bayer--Eisenbud \cite[\S8]{BE95}; it is used in this form in the proof of Proposition~\ref{prop:BEsection}.
\end{remark}

\section{The map \texorpdfstring{$\rho$}{rho}}\label{sec:rho}

This section is not used in the proof of Theorem~\ref{thm:alpha}. It records a second description of $\Psi$ through the extension classes $[V_{Y'}]$ of \eqref{eq:EY}: the morphism $\rho\colon\SC\to\PP(\Ext^1(\Omega_C,L_C))$, $Y'\mapsto[V_{Y'}]$, its local sections and its differential, and the identity $\beta\circ\Psi=\rho$, which gives a second proof of Proposition~\ref{prop:dPsi}.

\subsection{The map \texorpdfstring{$\rho$}{rho} and its differential}\label{subsec:rho}

\begin{theorem}\label{thm:rho}
Let $Q$ be as in \S\ref{subsec:sections}, $A:=\Ext^1(K_C,L_C)=\Ext^1(\Omega_C,L_C)\cong k^{g-2}$ and $\PP:=\PP(A)\cong\PP^{g-3}$.
\begin{enumerate}
\item There is a morphism $\rho\colon\SC\to\PP$ sending $Y'$ to the class of the extension
\[
0\to L_C\to V_{Y'}\to K_C\to0,\qquad V_{Y'}:=E_{Y'}\otimes L_C,
\]
where $E_{Y'}:=\pi'_*\OO_{Y'}(C)$, $\pi'\colon Y'\to C$ the ruling for which $C$ is a section, sits in $0\to\OO_C\to E_{Y'}\to Q\to0$. Moreover $Y'=\PP(E_{Y'})$ with $C$ the section defined by $E_{Y'}\to Q$, and $V_{Y'}\cong E_{Y'}^{\vee}\otimes K_C$.
\item Every point of $\PP$ has an open neighbourhood $U$ with a morphism $\Gamma_U\colon U\to\SC$ such that $\rho\circ\Gamma_U=\id_U$ and $\Gamma_U(\rho(Y'))\in\GH\cdot Y'$ for every $Y'\in\rho^{-1}(U)$.
\end{enumerate}
\end{theorem}

\begin{proof}
(1) Let $C\times\SC\hookrightarrow\mathcal{Y}\hookrightarrow\PM\times\SC$ be the universal scroll, as in Theorem~\ref{thm:Psi}. Let $p\colon\mathcal{Y}\to C\times\SC$ be the relative ruling for which $C\times\SC$ is a section and $\mathcal{V}:=p_*\OO_{\mathcal{Y}}(1)$; then $\mathcal{Y}=\PP(\mathcal{V})$ and
\[
\OO_{\mathcal{Y}}(1)=\OO_{\PP(\mathcal{V})}(1)=\OO_{\PP(\mathcal{V})}(C\times\SC),
\]
since $C\times\SC=\mathcal{Y}\cap(H\times\SC)$. Further $\OO_{C\times\SC}(1)=\OO_{C\times\SC}(C\times\SC)=Q\boxtimes\OO_{\SC}$. Now $C\times\SC$ is a section of $\PP(\mathcal{V})\to C\times\SC$ and hence comes with a surjection $\mathcal{V}\to Q\boxtimes\OO_{\SC}$; let $\mathcal{K}$ be its kernel. We claim $\mathcal{K}=\OO_{C\times\SC}$. From $0\to\OO_{\PP(\mathcal{V})}(-(C\times\SC))\to\OO_{\PP(\mathcal{V})}\to\OO_{C\times\SC}\to0$ we get $0\to\OO_{\PP(\mathcal{V})}\to\OO_{\PP(\mathcal{V})}(1)\to\OO_{C\times\SC}(1)\to0$, and pushing forward by $p$ ($R^1p_*\OO_{\PP(\mathcal{V})}=0$),
\begin{equation}\label{eq:ext}
0\to\OO_{C\times\SC}\to\mathcal{V}\to Q\boxtimes\OO_{\SC}\to0 .
\end{equation}
Tensoring with $L_C\boxtimes\OO_{\SC}$ we have
\begin{equation}\label{eq:V}
0\to L_C\boxtimes\OO_{\SC}\to\mathcal{V}\otimes(L_C\boxtimes\OO_{\SC})\to K_C\boxtimes\OO_{\SC}\to0 .
\end{equation}
On $C\times A$ there is a universal extension $0\to L_C\boxtimes\OO_A\to\mathcal{V}_A\to K_C\boxtimes\OO_A\to0$, of class $\id_A\in A\otimes A^{\vee}\subset A\otimes H^0(\OO_A)=\Ext^1_{C\times A}(K_C\boxtimes\OO_A,L_C\boxtimes\OO_A)$, and every extension of $K_C\boxtimes\OO_T$ by $L_C\boxtimes\OO_T$ over a scheme $T$ is its pull-back along a unique morphism $T\to A$ (base change, using $\Hom(K_C,L_C)=0$). Hence \eqref{eq:V} is the pull-back of $\mathcal{V}_A$ along a morphism $\nu\colon\SC\to A$. The fibre of \eqref{eq:ext} at $Y'$ is $0\to\OO_C\to E_{Y'}\to Q\to0$ with $E_{Y'}=\mathcal{V}|_{C\times\{Y'\}}=\pi'_*\OO_{Y'}(C)$, and it is nowhere split since $\SC$ consists of smooth scrolls: $E_{Y'}\cong\OO(a_1')\oplus\OO(a_2')$ with $a_i'\ge1$, so $E_{Y'}\not\cong\OO_C\oplus Q$. Hence $\nu$ takes values in $A\smallsetminus\{0\}$, and $\rho:=[\nu]\colon\SC\to A\smallsetminus\{0\}\to\PP$ is the required morphism. Finally $Y'=\PP(E_{Y'})$ with $C$ the section defined by $E_{Y'}\to Q$, and $E_{Y'}\otimes L_C=E_{Y'}\otimes Q^{\vee}\otimes K_C\cong E_{Y'}^{\vee}\otimes K_C$ because $\det E_{Y'}=Q$.

(2) Consider the universal extension $0\to L_C\boxtimes\OO_{\PP}(1)\to\mathcal{V}_{\PP}\to K_C\boxtimes\OO_{\PP}\to0$ on $C\times\PP$. Tensoring with $L_C^{-1}\boxtimes\OO_{\PP}(-1)$ we get
\[
0\to\OO_{C\times\PP}\to\mathcal{V}'\to Q\boxtimes\OO_{\PP}(-1)\to0,\qquad \mathcal{V}':=\mathcal{V}_{\PP}\otimes(L_C^{-1}\boxtimes\OO_{\PP}(-1)).
\]
The surjection $\mathcal{V}'\to Q\boxtimes\OO_{\PP}(-1)$ gives a section $i\colon C\times\PP\to\PP(\mathcal{V}')$ of $p\colon\PP(\mathcal{V}')\to C\times\PP$ such that $i^*\OO_{\PP(\mathcal{V}')}(1)=Q\boxtimes\OO_{\PP}(-1)$. Let $\mathrm{pr}\colon C\times\PP\to\PP$ and $\mathcal{W}:=\mathrm{pr}_*p_*\OO_{\PP(\mathcal{V}')}(1)=\mathrm{pr}_*\mathcal{V}'$. For $v\in\PP$, $\mathcal{W}_v=H^0(\mathcal{V}'|_{C\times\{v\}})=H^0(E_v)$ with $E_v\cong\OO(a_1')\oplus\OO(a_2')$, $a_i'\ge1$, $a_1'+a_2'=g-1$; so $h^0(E_v)=g+1=M+1$ and $h^1(E_v)=0$, and $\mathcal{W}$ is a vector bundle of rank $M+1$. Embed $\PP(\mathcal{V}')$ in $\PP(\mathcal{W})$ over $\PP$ by $|\OO_{\PP(\mathcal{V}')}(1)|$. Pushing forward the sequence above,
\[
0\to\OO_{\PP}\to\mathcal{W}\to H^0(Q)\otimes\OO_{\PP}(-1)\to0,
\]
and the image of $i(C\times\PP)$ lies in $\PP\big(H^0(Q)\otimes\OO_{\PP}(-1)\big)=\PP(H^0(Q))\times\PP=H\times\PP$, where it is $C\times\PP$, because $i(C\times\{v\})$ is embedded by $|i^*\OO(1)|=|\OO_C(1)|$. Given $v\in\PP$, let $U\ni v$ be an open set on which $\OO_{\PP}(-1)$ is trivial, e.g.\ the complement of a hyperplane, and choose a splitting of the sequence over $U$ (it exists, as $\Ext^1(H^0(Q)\otimes\OO_{\PP}(-1),\OO_{\PP})=H^0(Q)^{\vee}\otimes H^1(\OO_{\PP}(1))=0$). This gives $\mathcal{W}|_U\cong(k\oplus H^0(Q))\otimes\OO_U=H^0(\OO_{\PM}(1))\otimes\OO_U$ compatible with the quotients onto $H^0(Q)=H^0(\OO_H(1))$, hence $\PP(\mathcal{W})|_U\cong\PM\times U$ carrying $H\times U$ to $H\times U$ and $C\times U$ to $C\times U$. The composite $\PP(\mathcal{V}')|_U\hookrightarrow\PP(\mathcal{W})|_U\cong\PM\times U$ is a family of scrolls containing $C\times U$, i.e.\ a morphism $\Gamma_U\colon U\to\SC$. For $v\in U$, $\pi_*\OO_{\Gamma_U(v)}(C)=\mathcal{V}'_v=E_v$ with its extension, so $\rho\circ\Gamma_U=\id_U$. Finally, for $Y'\in\rho^{-1}(U)$, $\Gamma_U(\rho(Y'))$ and $Y'$ are both embeddings of $(\PP(E_{Y'}),C)$ by the complete linear series of $\OO(C)$ restricting to the given $C\subset H$; two such embeddings differ by an element of $\GH$, so $\Gamma_U(\rho(Y'))\in\GH\cdot Y'$.
\end{proof}

\begin{proposition}\label{prop:drho}
Let
\[
d\rho\colon H^0(\NY(-H))=T_{[Y]}\SC\longrightarrow T_{\rho(Y)}\PP=\Ext^1(\Omega_C,L_C)/k\,[V_Y]
\]
be the differential of $\rho$ at $[Y]$.
\begin{enumerate}
\item $d\rho$ is surjective.
\item $\ker d\rho=\Im c=T_{[Y]}(\GH\cdot Y)$.
\end{enumerate}
\end{proposition}

\begin{proof}
(1) Choose an open $U\subset\PP$ with $\rho(Y)\in U$ and a section $\Gamma_U\colon U\to\SC$ of $\rho$ as in Theorem~\ref{thm:rho}(2). Then $\rho\circ\Gamma_U=\id_U$, and differentiating at $\rho(Y)$,
\[
d\rho_{\Gamma_U(\rho(Y))}\circ d\Gamma_U=\id\colon T_{\rho(Y)}\PP\to T_{\rho(Y)}\PP .
\]
By Theorem~\ref{thm:rho}(2), $\Gamma_U(\rho(Y))=\gamma Y$ for some $\gamma\in\GH$, so $d\rho_{\gamma Y}$ is surjective.

\emph{Claim.} $\GH$ acts on $\SC$ and $\rho$ is $\GH$-invariant: $\rho\circ\gamma=\rho$ for $\gamma\in\GH$.

Granting the claim, $d\rho_{\gamma Y}\circ d\gamma_Y=d\rho_Y$, and $d\gamma_Y\colon T_Y\SC\to T_{\gamma Y}\SC$ is an isomorphism since $\gamma$ is an automorphism of $\SC$; hence $\operatorname{rank}d\rho_Y=\operatorname{rank}d\rho_{\gamma Y}$, and $d\rho_Y$ is surjective.

\emph{Proof of the claim.} Let $Y'\in\SC$ and $\gamma\in\GH$. Since $\gamma$ is a projectivity with $\gamma|_H=\id$, it maps $Y'$ onto a scroll $\gamma Y'$ containing $C$, so $\GH$ acts on $\SC$. Moreover $\gamma$ maps lines to lines and fixes $C$ pointwise, hence carries the fibres of $\pi'\colon Y'\to C$ (the lines of $Y'$ meeting $C$ in one point) to those of the ruling of $\gamma Y'$, and $\pi'(\gamma y)=\pi'(y)$; that is, the square
\[
\begin{tikzcd}[row sep=small]
Y'\arrow[r,"\gamma","\sim"']\arrow[d,"\pi'"'] & \gamma Y'\arrow[d,"\pi'"]\\
C\arrow[r,"\id"] & C
\end{tikzcd}
\]
commutes, with $\gamma|_C=\id_C$. Hence $\gamma^*$ induces an isomorphism $\pi'_*\OO_{\gamma Y'}(C)\xrightarrow{\sim}\pi'_*\OO_{Y'}(C)$ compatible with the restriction maps to $\OO_C(C)=Q$, i.e.\ an isomorphism $E_{\gamma Y'}\cong E_{Y'}$ inducing the identity on $Q$. It then carries the kernel $\OO_C$ of $E_{\gamma Y'}\to Q$ onto the kernel $\OO_C$ of $E_{Y'}\to Q$, necessarily by a nonzero scalar, so the extension classes of $E_{\gamma Y'}$ and $E_{Y'}$ in $\Ext^1(Q,\OO_C)$ are proportional. Therefore $\rho(\gamma Y')=\rho(Y')$ in $\PP$, and the claim is proved.

(2) By the claim, the two morphisms $\GH\times\SC\to\PP$, $(\gamma,Y')\mapsto\rho(\gamma Y')$ and $(\gamma,Y')\mapsto\rho(Y')$, agree on closed points; since $\GH\times\SC$ is reduced near $\GH\times\{Y\}$ (Lemma~\ref{lem:SC}(2)), they agree as morphisms there. Differentiating at $(1,Y)$ along $\GH$ gives $d\rho\circ c=0$, i.e.\ $\Im c\subset\ker d\rho$. By (1) and Lemma~\ref{lem:SC}(3),
\[
\dim\ker d\rho=h^0(\NY(-H))-\dim\PP=(M+1)+(2b-e-2)-(2b-e-2)=M+1=\dim\Im c ,
\]
so $\ker d\rho=\Im c$.
\end{proof}

\subsection{\texorpdfstring{$\Psi=\rho$}{Psi = rho}}\label{subsec:Psirho}

\begin{proposition}\label{prop:Psi=rho}
For every $Y'\in\SC$ one has $\beta(\Psi(Y'))=\rho(Y')$ in $\PP(\Ext^1(\Omega_C,L_C))$.
\end{proposition}

\begin{proof}
$\rho(Y')=[V_{Y'}]$ by Theorem~\ref{thm:rho}(1), and $\beta(\Psi(Y'))\in k^{\times}[V_{Y'}]$ by Proposition~\ref{prop:BEsection}(2) (or, more precisely, $\beta(\Psi(Y'))=[\Ct']=\yt'_h[V_{Y'}]$ with $\yt'_h\ne0$ by Proposition~\ref{prop:section}).
\end{proof}

\begin{remark}[Second proof of Proposition~\ref{prop:dPsi}]
By Proposition~\ref{prop:Psi=rho}, the two morphisms $\beta\circ\Psi$ and $\rho$ from $\SC$ to $\PP(\Ext^1(\Omega_C,L_C))$ agree on closed points ($\beta$ being the linear isomorphism of Lemma~\ref{lem:beta}). Since $\SC$ is reduced near $[Y]$ (Lemma~\ref{lem:SC}(2)), they agree as morphisms near $[Y]$, hence
\[
\bar\beta\circ d\Psi=d\rho\quad\text{at }[Y],
\]
where $\bar\beta\colon H^0(N_C\otimes L_C)/k\,s_C\xrightarrow{\sim}\Ext^1(\Omega_C,L_C)/k[V_Y]$ is induced by $\beta$ (note $\beta(s_C)=[\Ct]=\yt_h[V_Y]$). By Proposition~\ref{prop:drho}, $d\rho$ is surjective with kernel $\Im c$; therefore $d\Psi$ is surjective with kernel $\Im c$. This is Proposition~\ref{prop:dPsi}.
\end{remark}

\section{The class of the hyperplane section of the carpet}\label{sec:direct}

This section gives a proof of Proposition~\ref{prop:BEsection} which does not use Theorem~\ref{thm:extension}: it computes directly, for every $Y'\in\SC$, the class in $\Ext^1(\Omega_C,L_C)$ of the ribbon $\Yt'\cap H$. Parts (1), (2), (3) of the following proposition are parts (2), (1), (3) of Proposition~\ref{prop:BEsection}.

\begin{proposition}\label{prop:section}
Let $Y'\in\SC$, let $\yt'\in\Ext^1(\Omega_{Y'},K_{Y'})=H^1(T_{Y'}\otimes K_{Y'})$ be the class of the carpet $\Yt'$ (the restricted cotangent sequence $0\to K_{Y'}\to\Omega_{\Yt'}|_{Y'}\to\Omega_{Y'}\to0$; one has $\yt'=\partial(s')$), and let $\yt'_h\in k$ be its horizontal component, i.e.\ its image in $\Ext^1(p^*\Omega_{\PP^1},K_{Y'})=H^1(\Omega_{Y'/\PP^1})=k$ with respect to the basis given by the relative Euler class. Then the class of the restricted cotangent sequence of the ribbon $\Ct'=\Yt'\cap H$ is
\[
[\Ct']=\yt'_h\cdot[V_{Y'}]\in\Ext^1(\Omega_C,L_C),\qquad\text{and}\qquad \yt'_h\neq 0 .
\]
Consequently:
\begin{enumerate}
\item $\Yt'\cap H=\Ct_t$ for the $t\in H^0(N_C\otimes L_C)$ with $\beta(t)\in k^{\times}[V_{Y'}]$; i.e.\ $\Psi(Y')=[\beta^{-1}[V_{Y'}]]$.
\item $\Psi\circ\Gamma_U=\id_U$ for every $\Gamma_U$ as in Theorem~\ref{thm:Gamma}(2).
\item $\Gamma_U(\Psi(Y'))\in\GH\cdot Y'$ for every $Y'\in\Psi^{-1}(U)$.
\end{enumerate}
\end{proposition}

\begin{proof}
Write $Y$ for $Y'$, $p$ for the ruling, $\yt$, $\Ct_H:=\Yt\cap H$. Let $q\colon\Omega_Y|_C\to\Omega_C$ be the quotient map (kernel $N_{C/Y}^{\vee}=Q^{\vee}$) and $\sigma\colon\Omega_C\to\Omega_Y|_C$ the splitting induced by $p$ via $C\cong\PP^1$, so $q\circ\sigma=\id$.

\emph{Step (i): $\yt|_C=q^*[\Ct_H]$.} Pulling back differentials along $\Ct_H\hookrightarrow\Yt$ and restricting to $C$ gives a map $\Omega_{\Yt}|_C\to\Omega_{\Ct_H}|_C$ which on the subsheaves is the identification $(I_Y/I_{\Yt})|_C=I_C/I_{\Ct_H}$ ($=L_C$, Lemma~\ref{lem:RC}) and on the quotients is $q$. A morphism of extensions which is the identity on the left term and $q$ on the right term exhibits the source as the pull-back of the target along $q$.

\emph{Step (ii): $[\Ct_H]=\sigma^*(\yt|_C)$.} From $0\to Q^{\vee}\to\Omega_Y|_C\to\Omega_C\to0$ and $\Hom(-,L_C)$, the kernel of $q^*\colon\Ext^1(\Omega_C,L_C)\to\Ext^1(\Omega_Y|_C,L_C)$ is the image of $\Hom(Q^{\vee},L_C)=H^0(Q\otimes L_C)=H^0(\OO_{\PP^1}(-2))=0$, so $q^*$ is injective and $[\Ct_H]$ is the unique class with $q^*[\Ct_H]=\yt|_C$. Hence $[\Ct_H]=\sigma^*q^*[\Ct_H]=\sigma^*(\yt|_C)$.

\emph{Step (iii): factorisation.} The map $\yt\mapsto\sigma^*(\yt|_C)$ is ``pull back along $p^*\Omega_{\PP^1}\to\Omega_Y$, then restrict to $C$'' (both are pull-back along the composite $\Omega_C=(p^*\Omega_{\PP^1})|_C\to\Omega_Y|_C$). It therefore factors through
\[
\Ext^1(\Omega_Y,K_Y)\xrightarrow{\ r_h\ }\Ext^1(p^*\Omega_{\PP^1},K_Y)=H^1(K_Y\otimes p^*T_{\PP^1})=H^1(\Omega_{Y/\PP^1})\xrightarrow{\ \operatorname{res}_C\ }\Ext^1(\Omega_C,L_C),
\]
using $K_Y\otimes p^*T_{\PP^1}=K_{Y/\PP^1}=\Omega_{Y/\PP^1}$. Now $H^1(\Omega_{Y/\PP^1})=H^0(R^1p_*\Omega_{Y/\PP^1})=k$ is spanned by the class $u$ of the relative Euler sequence $0\to\Omega_{Y/\PP^1}\to p^*E_Y\otimes\OO_{\PP(E_Y)}(-1)\to\OO_Y\to0$ (tensored with $p^*\Omega_{\PP^1}$ to view it in $\Ext^1(p^*\Omega_{\PP^1},K_Y)$). Restricting this sequence to $C$, where $\OO_{\PP(E_Y)}(-1)|_C=Q^{\vee}$ and $\Omega_{Y/\PP^1}|_C=Q^{\vee}$, and tensoring with $\Omega_C=K_C$ gives
\[
0\to Q^{\vee}\otimes K_C=L_C\to E_Y\otimes Q^{\vee}\otimes K_C=V_Y\to K_C\to0,
\]
i.e.\ $\operatorname{res}_C(u)=[V_Y]$ (this is \cite[Claim~7.6]{DM25}). Writing $r_h(\yt)=\yt_h\cdot u$ we get $[\Ct_H]=\yt_h\,[V_Y]$.

\emph{Step (iv): $\yt_h\ne0$.} Since $p_*(K_Y\otimes p^*T_{\PP^1})=T_{\PP^1}\otimes p_*K_Y=0$, the Leray map identifies $H^1(K_Y\otimes p^*T_{\PP^1})$ with $H^0(R^1p_*(K_Y\otimes p^*T_{\PP^1}))=H^0(R^1p_*(T_Y\otimes K_Y))$ (using $R^1p_*(T_{Y/\PP^1}\otimes K_Y)=R^1p_*(p^*K_{\PP^1})=0$), and under this identification $r_h$ is the Leray map $H^1(T_Y\otimes K_Y)\to H^0(R^1p_*(T_Y\otimes K_Y))$. By Lemma~\ref{lem:pushforward} the image of $\yt=\partial(s)$ under this map is nonzero.

(1) $\beta(\Psi(Y'))=[\Ct']=\yt'_h[V_{Y'}]$ with $\yt'_h\ne0$, and $\beta$ is injective.

(2) Let $u=[t]\in U$ and $Y':=\Gamma_U(u)$. Then $[V_{Y'}]\in k^{\times}\beta(t)$ by Theorem~\ref{thm:Gamma}(2), so $\beta(\Psi(Y'))\in k^{\times}\beta(t)$ by (1), and $\Psi(Y')=[t]=u$.

(3) Let $Y'\in\Psi^{-1}(U)$, $\Psi(Y')=[t]$ and $Y'':=\Gamma_U([t])$. By (1) and Theorem~\ref{thm:Gamma}(2), $[V_{Y'}]$ and $[V_{Y''}]$ both lie in $k^{\times}\beta(t)$; so the extensions \eqref{eq:EY} of $Y'$ and $Y''$ have proportional classes, and there is an isomorphism $E_{Y'}\cong E_{Y''}$ inducing the identity on $Q$ (and a nonzero scalar on $\OO_C$). Hence $(\PP(E_{Y'}),C)\cong(\PP(E_{Y''}),C)$ over $C$, and $Y'$, $Y''$ are two embeddings of this pair by the complete linear series $|\OO(C)|$ which agree on $C$. They differ by a $\gamma\in PGL_{M+1}(k)$ with $\gamma|_C=\id$, hence $\gamma|_H=\id$ since $C$ spans $H$; so $Y''=\gamma Y'$ with $\gamma\in\GH$.
\end{proof}


Conversely, Proposition~\ref{prop:section}(1) applied to $Y'=\Gamma_U(u)$, $u=[\Ct]$, recovers Theorem~\ref{thm:extension}(3) in the form $\Yt_{\Ct}\cap H=\Ct$.

\appendix

\section{The compatibility \texorpdfstring{$h\circ f=2g$}{h o f = 2g}}\label{app:hfg}

We keep the notation of the proof of Theorem~\ref{thm:alpha}: $\Hc(-H)\to\NYt(-H)$ is the first map of \eqref{eq:3.2}$(-H)$, denoted $\kappa_0$ on global sections, $F$ is the connecting map of \eqref{eq:3.65}, $f=F\circ\kappa_0$, $g$ is the connecting map of \eqref{eq:3.3}$(-H)$, and $h=H^1$ of the last map of \eqref{eq:3.2}$(-H+K_Y)$, i.e.\ restriction of homomorphisms to $I_Y^2/I_YI_{\Yt}\cong K_Y^{\otimes 2}$. We identify $\sHom(K_Y,\OO_Y(-H))\cong\sHom(K_Y^{\otimes2},K_Y(-H))$ via $\vartheta\mapsto\vartheta\otimes\id_{K_Y}$.

\begin{lemma}\label{lem:hfg}
$h\circ f=2\,g$. Consequently $\Im(h\circ f)=\Im g$, and the right square of diagram \eqref{diag:a} commutes with $2\cdot\id$ as its right-hand vertical map.
\end{lemma}

\begin{proof}
We use \v{C}ech cocycles on an affine cover $\{U_i\}$ of $\Yt$ (equivalently of $Y$), $U_{ij}=U_i\cap U_j$, and write $L:=I_Y/I_{\Yt}\cong K_Y$. Fix $\varphi\in H^0(\Hc(-H))$, i.e.\ $\varphi\colon I_{\Yt}/I_Y^2\to\OO_Y(-H)$, and let $\bar\varphi=\kappa_0(\varphi)=\varphi\circ\pi$ with $\pi\colon I_{\Yt}/I_YI_{\Yt}\to I_{\Yt}/I_Y^2$; note $\bar\varphi(I_Y^2)=0$.

\emph{Cocycle for $f(\varphi)=F(\bar\varphi)$.} Choose $\OO_{\Yt}$-linear lifts $\Phi_i\colon I_{\Yt}/I_{\Yt}^2|_{U_i}\to\OO_{\Yt}(-\Ht)|_{U_i}$ of $\bar\varphi$. Then $\Phi_{ij}=\Phi_i-\Phi_j$ takes values in $L(-H)$, and $f(\varphi)=[\{\Phi_{ij}\}]$. Since $h$ is restriction along $L^{\otimes2}\cong I_Y^2/I_YI_{\Yt}$,
\[
h(f(\varphi))=\big[\{\Phi_{ij}|_{L^{\otimes 2}}\}\big],\qquad \Phi_{ij}|_{L^{\otimes2}}(\bar x\otimes\bar y)=\Phi_i(xy)-\Phi_j(xy)\quad(x,y\in I_Y).
\]

\emph{Cocycle for $g(\varphi)$.} Choose $\OO_Y$-linear lifts $\psi_i\colon I_Y/I_Y^2|_{U_i}\to\OO_Y(-H)$ with $\psi_i|_{I_{\Yt}/I_Y^2}=\varphi$; then $\psi_{ij}=\psi_i-\psi_j$ factors through $L$ and $g(\varphi)=[\{\psi_{ij}\}]$.

\emph{Comparison.} Put $B_i(x,y):=\Phi_i(xy)\in L(-H)$ for $x,y\in I_Y$ (well defined since $\bar\varphi(I_Y^2)=0$). Then: (i) $B_i$ is symmetric, $\OO$-bilinear, and vanishes if $x$ or $y$ lies in $I_Y^2$ (if $x=\sum x'_kx''_k$, then $\Phi_i(x'_kx''_ky)=\bar x''_k\cdot\Phi_i(x'_ky)\in L\cdot L(-H)=0$), so it descends to $I_Y/I_Y^2$; (ii) if $x\in I_{\Yt}$ then $B_i(x,y)=\bar y\cdot\Phi_i(x)=\varphi(\bar x)\otimes\bar y$, because multiplication by $\bar y\in L$ on $\OO_{\Yt}(-\Ht)$ factors through $\OO_Y(-H)$; symmetrically for $y\in I_{\Yt}$. Put $B'_i(x,y):=\psi_i(x)\otimes\bar y+\psi_i(y)\otimes\bar x$. By (ii), $B_i=B'_i$ whenever one argument lies in $I_{\Yt}/I_Y^2$, so $C_i:=B_i-B'_i$ descends to a section of $\sHom(L^{\otimes2},L(-H))=L^{-1}(-H)$ over $U_i$. Over $U_{ij}$, as forms on $L\times L$,
\[
B_i-B_j=(C_i-C_j)+(B'_i-B'_j),\qquad (B'_i-B'_j)(u,v)=\psi_{ij}(u)\otimes v+\psi_{ij}(v)\otimes u=2\,\psi_{ij}(u)\otimes v,
\]
the last step because $L$ is a line bundle ($\vartheta(u)\otimes v=\vartheta(v)\otimes u$). Hence $\{B_i-B_j\}=\delta\{C_i\}+2\{\psi_{ij}\otimes\id_L\}$, i.e.\ $h(f(\varphi))=2g(\varphi)$.
\end{proof}

\begin{thebibliography}{GGP08}

\bibitem[AHL10]{AHL10} M.~Artebani, J.~Hausen and A.~Laface, \emph{On Cox rings of K3 surfaces}, Compos.\ Math.\ \textbf{146} (2010), no.~4, 964--998.

\bibitem[BE95]{BE95} D.~Bayer, D.~Eisenbud, \emph{Ribbons and their canonical embeddings}, Trans.\ Amer.\ Math.\ Soc.\ \textbf{347} (1995), 719--756.

\bibitem[BM24]{BM24} P.~Bangere, J.~Mukherjee, \emph{Extendability of projective varieties via degeneration to ribbons with applications to Calabi--Yau threefolds}, arXiv:2409.03960.

%\bibitem[BM25]{BM25} P.~Bangere, J.~Mukherjee, \emph{Extendability of K3 surfaces without Gaussian maps and classification of prime Fano threefolds}, draft. %% adjust: this is the manuscript whose \S3 is reproduced in \S\ref{subsec:reduction}

\bibitem[BMR21]{BMR21}
Bangere, Purnaprajna; Mukherjee, Jayan; Raychaudhury, Debaditya
\emph{K3 carpets on minimal rational surfaces and their smoothings.}
Internat. J. Math.32(2021), no.6, Paper No. 2150032, 20 pp.

\bibitem[AHL10]{AHL10} M.~Artebani, J.~Hausen and A.~Laface, \emph{On Cox rings of K3 surfaces}, Compos.\ Math.\ \textbf{146} (2010), no.~4, 964--998.

\bibitem[But94]{But94} D.~C.~Butler, \emph{Normal generation of vector bundles over a curve}, J.\ Differential Geom.\ \textbf{39} (1994), 1--34.

\bibitem[CLM93]{CLM93} C.~Ciliberto, A.~F.~Lopez, R.~Miranda, \emph{Projective degenerations of K3 surfaces, Gaussian maps, and Fano threefolds}, Invent.\ Math.\ \textbf{114} (1993), 641--667.

\bibitem[CLM98]{CLM98} C.~Ciliberto, A.~F.~Lopez, R.~Miranda, \emph{Classification of varieties with canonical curve section via Gaussian maps on canonical curves}, Amer.\ J.\ Math.\ \textbf{120} (1998), 1--21.

\bibitem[DM24]{DM24} A.~Deopurkar, J.~Mukherjee, \emph{Syzygies of canonical ribbons on higher genus curves}, arXiv:2412.05500.

\bibitem[DM25]{DM25} A.~Deopurkar, J.~Mukherjee, \emph{Stable reduction of ribbons and K-trivial carpets on surface scrolls}, preprint.

\bibitem[ES18]{ES18} D.~Eisenbud, F.-O.~Schreyer, \emph{Equations and syzygies of K3 carpets and unions of scrolls}, arXiv:1804.08011.

\bibitem[GGP08]{GGP08} F.~J.~Gallego, M.~Gonz\'alez, B.~P.~Purnaprajna, \emph{K3 double structures on Enriques surfaces and their smoothings}, J.\ Pure Appl.\ Algebra \textbf{212} (2008), 981--993.

\bibitem[GP97]{GP97} F.~J.~Gallego, B.~P.~Purnaprajna, \emph{Degenerations of K3 surfaces in projective space}, Trans.\ Amer.\ Math.\ Soc.\ \textbf{349} (1997), 2477--2492.

\bibitem[Huy16]{Huy16} D.~Huybrechts, \emph{Lectures on K3 surfaces}, Cambridge Studies in Advanced Mathematics \textbf{158}, Cambridge University Press, Cambridge, 2016.

\bibitem[Isk77]{Isk77} V.~A.~Iskovskikh, \emph{Fano threefolds. I}, Izv.\ Akad.\ Nauk SSSR Ser.\ Mat.\ \textbf{41} (1977), 516--562.

\bibitem[Isk78]{Isk78} V.~A.~Iskovskikh, \emph{Fano threefolds. II}, Izv.\ Akad.\ Nauk SSSR Ser.\ Mat.\ \textbf{42} (1978), 506--549.

\bibitem[Kle81]{Kle81} J.~O.~Kleppe, \emph{The Hilbert-flag scheme, its properties and its connection with the Hilbert scheme. Applications to curves in 3-space}, preprint (part of thesis), University of Oslo, 1981.

\bibitem[Lvo92]{Lvo92} S.~L'vovsky, \emph{Extensions of projective varieties and deformations. I, II}, Michigan Math.\ J.\ \textbf{39} (1992), 41--51, 65--70.

\bibitem[Mat89]{Mat89} H.~Matsumura, \emph{Commutative ring theory}, Cambridge Studies in Advanced Mathematics \textbf{8}, Cambridge University Press, Cambridge, 1989.

\bibitem[Muk88]{Muk88} S.~Mukai, \emph{Curves, K3 surfaces and Fano 3-folds of genus $\le 10$}, in: Algebraic geometry and commutative algebra, Vol.~I, Kinokuniya, Tokyo, 1988, 357--377.

\bibitem[SD74]{SD74} B.~Saint-Donat, \emph{Projective models of K-3 surfaces}, Amer.\ J.\ Math.\ \textbf{96} (1974), 602--639.

\bibitem[Ser06]{Ser06} E.~Sernesi, \emph{Deformations of algebraic schemes}, Grundlehren der Mathematischen Wissenschaften \textbf{334}, Springer, Berlin, 2006.

\bibitem[Zak91]{Zak91} F.~L.~Zak, \emph{Some properties of dual varieties and their applications in projective geometry}, in: Algebraic geometry (Chicago, IL, 1989), Lecture Notes in Math.\ \textbf{1479}, Springer, Berlin, 1991, 273--280.


\end{thebibliography}

\end{document}
